% !TeX program = pdflatex
% Integral second-order Apery-like recurrences and rational elliptic surfaces.
% Self-contained; builds with pdflatex; bibliography inline.
\documentclass[11pt,reqno]{amsart}

\usepackage[T1]{fontenc}
\usepackage[utf8]{inputenc}
\usepackage{amsmath,amssymb,amsthm}
\usepackage{mathtools}
\usepackage{array}
\usepackage{enumitem}
\usepackage[margin=0.95in]{geometry}
\usepackage{microtype}
\usepackage[hidelinks]{hyperref}

\theoremstyle{plain}
\newtheorem{theorem}{Theorem}[section]
\newtheorem{proposition}[theorem]{Proposition}
\newtheorem{lemma}[theorem]{Lemma}
\newtheorem{corollary}[theorem]{Corollary}
\theoremstyle{definition}
\newtheorem{definition}[theorem]{Definition}
\theoremstyle{remark}
\newtheorem{remark}[theorem]{Remark}

\setlength{\emergencystretch}{3em}

% Modest vertical tightening (amsart 11pt is generous by default).
\setlength{\abovedisplayskip}{5pt plus 2pt minus 2pt}
\setlength{\belowdisplayskip}{5pt plus 2pt minus 2pt}
\setlength{\abovedisplayshortskip}{3pt plus 2pt}
\setlength{\belowdisplayshortskip}{3pt plus 2pt}
\makeatletter
\def\thm@space@setup{\thm@preskip=5pt \thm@postskip=5pt}
\makeatother

% Named main theorems: A (conditional classification), B (unconditional
% necessary conditions), C (the scan).
\theoremstyle{plain}
\newtheorem{namedthm}{Theorem}
\renewcommand{\thenamedthm}{\Alph{namedthm}}
\newtheorem{namedcor}[namedthm]{Corollary}

\newcommand{\Z}{\mathbb{Z}}
\newcommand{\Q}{\mathbb{Q}}
\newcommand{\R}{\mathbb{R}}
\newcommand{\PP}{\mathbb{P}}
\newcommand{\HH}{\mathbb{H}}
\DeclareMathOperator{\PSL}{PSL}
\DeclareMathOperator{\Sym}{Sym}
\DeclareMathOperator{\ord}{ord}
\DeclareMathOperator{\lc}{lc}
\DeclareMathOperator{\Res}{Res}
\newcommand{\tht}{\theta_t}
\newcommand{\thq}{\theta_q}

% ---------------------------------------------------------------------------
% Evidence-class markers.  These are load-bearing and are never upgraded.
% ---------------------------------------------------------------------------
\newcommand{\Proved}{\textup{\textsf{[PROVED]}}}
\newcommand{\VerifiedR}[1]{\textup{\textsf{[VERIFIED: #1]}}}
\newcommand{\Numerical}[1]{\textup{\textsf{[NUMERICAL: #1]}}}
\newcommand{\Open}{\textup{\textsf{[OPEN]}}}

\begin{document}

\title[Integral second-order rows and elliptic surfaces]
      {Integral second-order Ap\'ery-like recurrences\\ and rational elliptic surfaces}
\author{Claude (Fable), with River}
\date{\today}

\begin{abstract}
We study the integral solutions of three-term recurrences
$(n+1)^2u_{n+1}=P(n)u_n-Q(n)u_{n-1}$, $\deg P,\deg Q\le2$, in two tiers.
\emph{Theorem A} is a conditional classification and is proved: if the
Picard--Fuchs operator of such a row is that of an elliptic surface over
$\PP^1_t$, then for the Zagier shape $(an^2+an+b,\;cn^2)$ all four singular
fibres are semistable, $\deg\mathcal J=12$ exactly --- by Riemann--Hurwitz
together with a twist-parity constraint $12\mid\deg\mathcal J$ --- and the
configuration is, up to quadratic twist, one of Beauville's six; no hypothesis on
the Kodaira dimension is needed, since a twist raises $\chi$ (two branch points
give a $K3$) without changing $\mathcal J$, the local system, or the row. The row
is the period expansion at a cusp in the integral coordinate, and the complete
list of cusp placements is Zagier's seven rows
$\mathbf A$--$\mathbf G$, with $\mathbf A,\mathbf C,\mathbf F$ the three
placements of $I_6I_3I_2I_1$ and $\mathbf E,\mathbf G$ the two of
$I_1I_1I_2I_8$. The correspondence is Zagier's (2009, \S7) with Beauville's
theorem; new here are the mechanism --- an exponent dictionary, Kodaira
admissibility, a cross-ratio invariant, and a cusp-move theorem enumerating the
placements --- and the completeness of the list. We prove the analogue for the
root-row shape $(An^2+\tfrac A2n+B,\;C(2n-1)^2)$: there $\deg\mathcal J=6$,
three of Herfurtner's configurations survive up to twist, and exactly two
integral rows realise them. Riemann--Hurwitz alone does not force $\chi=1$ in
either shape --- in the root-row case a mixed $III/III^*$ elliptic $K3$ is inside
the bound --- and it is the twist parity, not the degree bound, that settles it.
Beukers' row, the square root of Ap\'ery's $\zeta(3)$ recurrence, is not among
them --- it lives on $X_0(6)/W_6$, over which the
universal curve does not descend. \emph{Theorem B} collects the unconditional
necessary conditions --- Kodaira admissibility of the exponent differences,
integrality rigidity (rational characteristic roots are integers, so the
irrationality score is $\le-k$ on every rational cusp-move orbit), and the
Riemann--Hurwitz bound --- which exclude the square roots of Cooper's
$s_{10},s_{18}$ and an infinite $\arctan$-Pad\'e family. The unconditional
classification is exactly Zagier's conjecture that any integral solution has a
modular parametrisation, i.e.\ Bombieri--Dwork in rank two with four singular
points; we make the reduction precise and give the evidence: exhaustive scans
($1.27\times10^6$ integrality hits over $26$ Kodaira-admissible classes;
$3.0\times10^6$ eta-quotient pairs), two new elliptic-surface rows on
$I_1I_7\,II\,II$ and $I_3\,III\,III\,III$ (the latter with Chowla--Selberg period
$\operatorname{Im}\xi=-\Gamma(1/4)^4/(2^{7/2}3^{11/4}\pi)$), and the fact that
Ap\'ery's $\zeta(2)$ row and Beukers' row are the only rows in the box with
positive score and denominator exponent $k\ge2$.
\end{abstract}

\maketitle

% ===========================================================================
\section{Introduction}\label{sec:intro}
% ===========================================================================

\subsection*{The question}
Fix the three-term recurrence
\begin{equation}\label{eq:R2gen}
(n+1)^2u_{n+1}=P(n)\,u_n-Q(n)\,u_{n-1},\qquad u_0=1,\quad u_{-1}=0,
\end{equation}
with $P=an^2+bn+c$ and $Q=dn^2+en+f$ in $\Q[n]$. Which choices of $(P,Q)$ give
$u_n\in\Z$ for all $n$?

The question is Ap\'ery's, in the form Zagier gave it. Ap\'ery's proof of the
irrationality of $\zeta(2)$ runs on
$(n+1)^2u_{n+1}=(11n^2+11n+3)u_n+n^2u_{n-1}$, whose two solutions are integral,
respectively integral after multiplication by $\operatorname{lcm}(1,\dots,n)^2$.
Zagier \cite{Zagier2009} searched the sub-family $P=an^2+an+b$, $Q=-cn^2$ by
computer and found, besides degenerate and hypergeometric solutions, exactly six
sporadic integral rows $\mathbf A,\dots,\mathbf F$ (Ap\'ery's is $\mathbf D$),
all modular: each is the $t$-expansion of a weight-one form against a Hauptmodul
on a genus-zero group, and the six groups are the six semistable configurations
of Beauville \cite{Beauville1982}. Almkvist--Zudilin
\cite{AlmkvistZudilin2006} carried out the analogous search one order up, Cooper
\cite{Cooper2012} added $s_7,s_{10},s_{18}$, and Beukers \cite{Beukers1987}
supplied the framework in which all of these are Eichler integrals --- and, in his
Theorem~3, an irrationality proof for a period attached to the \emph{square root}
of Ap\'ery's $\zeta(3)$ recurrence.

Zagier's own \S7 already contains the answer to the second half of this, and it
should be quoted before anything is claimed. He observes that Beauville's six
stable families give Picard--Fuchs equations of the form \eqref{eq:R2gen}, that
two of Beauville's families ($V$ and $VI$) are isogenous to two others ($II$ and
$I$) and give the same equations, that his hypergeometric solutions \#2, \#11,
\#14 are also on Beauville's list, and that ``each of Beauville's surfaces can
give rise to several different Picard--Fuchs equations, depending on how one
chooses the Hauptmodul $t(z)$ \dots\ we must put one fibre at $t=0$ and one at
$\infty$, giving $12$ choices''. He then gives the correspondence
\[
\begin{array}{lcccc}
\text{Beauville family:} & IV & I,VI & III & II,V\\
\text{Picard--Fuchs equation:} & \mathbf A,\mathbf C,\mathbf F
 & \mathbf B,\ \#14 & \mathbf D & \mathbf E,\mathbf G,\ \#11
\end{array}
\]
and notes $t_{\mathbf C}(z)=t_{\mathbf A}(z)/(1+t_{\mathbf A}(z))$. So the fact
that $\mathbf A,\mathbf C,\mathbf F$ are three placements on one surface, and
that $\mathbf E,\mathbf G$ are two on another, is \emph{Zagier's}, not ours.

What is not settled there is why the answer is finite, why six, and whether a
solution can exist outside the imposed normalisation. Zagier says as much: after
listing $12$ modular solutions he writes that a complete classification of the
\emph{modular} solutions ``probably can be done'', that the most optimistic guess
is that the $12$ already found are complete, and that the passage from ``modular''
to ``all'' is the content of

\begin{quote}
\textbf{Conjecture} (Zagier \cite[p.~6]{Zagier2009}). \emph{Any integral solution
of the differential equation \textup{(1)}, where $P(t)$ is a non-degenerate
quadratic polynomial, has a modular parametrization.}
\end{quote}

\subsection*{What this paper does}
We give a two-tier answer. The classification splits into a \emph{local}
question --- which shapes of $(P,Q)$ are possible --- and a \emph{global} one ---
which shapes actually carry an elliptic surface. Both are settled by proof,
conditionally on the geometric hypothesis; the removal of that hypothesis is
Zagier's conjecture, and \S\ref{sec:BD} makes the reduction precise.

\begin{namedthm}[conditional classification]\label{thmA}
\Proved{} \textup{(}given Beauville's and Herfurtner's classifications; the
explicit lists are exact finite computations, marked as such in
\S\ref{sec:thmA}\textup{)}.
\begin{enumerate}[label=\textup{(\arabic*)},leftmargin=2.6em]
\item \emph{Zagier shape.} Let $a,b,c\in\Z$ with $c\ne0$ and $a^2\ne4c$, and let
$u_n\in\Z$ solve $(n+1)^2u_{n+1}=(an^2+an+b)u_n-cn^2u_{n-1}$, $u_0=1$,
$u_{-1}=0$. Suppose the local system $\ker L$ is, up to a rational gauge factor
and a quadratic character twist, $R^1\pi_*\Q$ of an elliptic surface
$\pi:\mathcal E\to\PP^1_t$ with non-constant $\mathcal J$. Then \emph{(a)} all
four local exponent differences vanish, so all four singular fibres are
semistable (up to the twist); \emph{(b)} $\deg\mathcal J\le12$ by
Riemann--Hurwitz and $12\mid\deg\mathcal J$ by twist parity
(Lemma~\ref{lem:twistparity}), so $\deg\mathcal J=12$ exactly and $\mathcal J$ is
a Belyi map; the \emph{untwisted} member of the surface's quadratic-twist class
is then semistable with $\sum e=12$, i.e.\ a \emph{rational} elliptic surface,
and is one of Beauville's six. No hypothesis on the Kodaira dimension is
required: a quadratic twist raises $\chi$ (two branch points give a $K3$) without
changing $\mathcal J$, the local system, or the row; \emph{(c)} $t$ is the Hauptmodul
normalised $t=q_h+O(q_h^2)$ at the MUM cusp and $u_n$ is its $t$-expansion
coefficient sequence --- the row is the period expansion at a cusp in the integral
coordinate; and \emph{(d)} the cross-ratio $\mathcal I=a^2/c$ lies in
$\{0,3,\tfrac92,-\tfrac{49}8,\tfrac{100}9,\tfrac{289}{72},-121\}$, the
configuration together with the choice of MUM cusp and of the cusp at $\infty$
determines the operator and hence $b$, and the complete list of such placements
is the seven rows $\mathbf A$--$\mathbf G$ of Table~\ref{tab:placements}: three
placements of $I_6I_3I_2I_1$ ($\mathbf A,\mathbf C,\mathbf F$), two of
$I_1I_1I_2I_8$ ($\mathbf E,\mathbf G$), one each of $I_1I_1I_1I_9$
($\mathbf B$) and $I_1I_1I_5I_5$ ($\mathbf D$).
\item \emph{Root-row shape.} Let the row be
$(An^2+\tfrac A2n+B,\;C(2n-1)^2)$ with $A,B,C\in\Z$, integral after a $\lambda^n$
rescaling, and make the same geometric hypothesis. Then $\rho=\tfrac12$ and
$\delta_\infty=0$: the two finite fibres are of type $III$ or $III^*$ and the
fibre at $\infty$ is a cusp. Riemann--Hurwitz gives $\deg\mathcal J\le6$ and
twist parity gives $\deg\mathcal J\equiv6\pmod{12}$, so $\deg\mathcal J=6$; the
untwisted member of the twist class has fibres $I_n,I_m,III,III$ with $n+m=6$ and
$\sum e=12$, hence is \emph{rational}, and is one of Herfurtner's \#38
$I_1I_5\,III\,III$, \#39 $I_2I_4\,III\,III$, \#40 $I_3I_3\,III\,III$, with
$\mathcal I=\tfrac{484}{125},0,-12$. Riemann--Hurwitz alone does \emph{not} force
$\chi=1$ here --- the mixed configuration $III$ at $t_1$, $III^*$ at $t_2$, $I_n$
at $0$, $I_m^*$ at $\infty$ with $n+m=6$ has $\sum e=24$, an elliptic $K3$, and
$\deg\mathcal J=6$ sits inside the bound --- but that configuration is the
quadratic twist with branch set $\{t_2,\infty\}$ of the corresponding
\#38/\#39/\#40 one, so it carries the same $\mathcal J$, the same local system and
the same row \textup{(}\S\ref{sec:thmA}\textup{)}. Exactly two integral rows
realise them:
$\sqrt{\mathrm{AZ}(11,5,125)}=(88,10,500)$ on \#38 with cusp width $h=1$, and
$\sqrt{\mathrm{AZ}(9,3,-27)}=(72,6,-108)$ on \#40 with $h=3$; the value
$\mathcal I=0$ forces $A=0$ and yields no non-degenerate integral row. In
particular \emph{Beukers' row} $(136,10,4)$, with $\mathcal I=1156$, satisfies the
hypothesis for no elliptic surface: its projective monodromy is the
Atkin--Lehner quotient $\Gamma_0(6)+6$, and over $\PP^1_t=X_0(6)/W_6$ the
universal elliptic curve does not descend --- only its symmetric square does.
\end{enumerate}
\end{namedthm}

Part (1) is, in substance, Zagier \S7 plus Beauville, and we claim no priority
for the correspondence. What is new is the mechanism and the completeness: the
exponent dictionary and the fact that Zagier's normalisation is \emph{forced} by
semistability rather than assumed (Theorem~\ref{thm:norm}); Kodaira admissibility
(Theorem~\ref{thm:kodaira}); the cross-ratio invariant $\mathcal I=a^2/d$ as a
finite \emph{necessary} condition, which turns Zagier's ``$12$ choices'' into a
finite check (Theorem~\ref{thm:finite}); a cusp-move theorem enumerating the
placements as a group action, proving the placement list complete and showing
that $\mathbf A,\mathbf C,\mathbf F$ carry \emph{two} distinct extension classes
on the one surface (\S\ref{sec:cusp}); the root-row case, which is not in Zagier;
the Riemann--Hurwitz bound removing ``rational'' from the hypothesis
(Proposition~\ref{prop:K1}); the scans; and the non-congruence classification of
\S\ref{sec:noncong}. Relative to Movasati--Reiter \cite{MovasatiReiter}, who list
the Heun equations attached to Herfurtner's families, we add the recurrence and
integrality side, the cross-ratio criterion, the root rows, and the scans.

\begin{namedthm}[unconditional necessary conditions]\label{thmB}
\Proved{} Let \eqref{eq:R2gen} be integral with $d\ne0$, $a^2\ne4d$ and
$Q(n)\ne0$ for $n\ge1$. Then, with no geometric hypothesis:
\begin{enumerate}[label=\textup{(\roman*)},leftmargin=2.6em]
\item \emph{(normalisation)} if $a^2-4d\notin\Q^{\times2}$ then the two movable
exponents agree, $\rho_1=\rho_2=\rho$, and $b=(1-\rho)a$, $e=-2\rho d$;
\item \emph{(Kodaira admissibility)} if $\ker L$ is, up to rational gauge and
quadratic twist, that of an elliptic surface, then every local exponent
difference lies in $\{0,\tfrac12,\tfrac13,\tfrac23\}\bmod1$; contrapositively,
any row violating this comes from no elliptic surface, in any gauge;
\item \emph{(integrality rigidity)} $\lambda_1,\lambda_2$ are algebraic integers,
so $\lambda_1\in\Q$ iff $\lambda_1,\lambda_2\in\Z$ iff $a^2-4d$ is a square; in
that case all three gaps $|\lambda_1|,|\lambda_2|,|\lambda_1-\lambda_2|$ are
positive integers and hence $|\lambda_2|\ge1$ at \emph{every} cusp placement, so
$\operatorname{score}=\log(1/|\lambda_2|)-k\le-k\le-1$ throughout the rational
orbit. Consequently a positive-score row has irrational characteristic roots;
\item \emph{(Riemann--Hurwitz)} for any elliptic surface of any Kodaira dimension
whose four singular fibres sit over the row's singular points,
$\deg\mathcal J\le6c+4e_3+3e_2-12\le12$, where $c,e_2,e_3$ count the points whose
projective local monodromy is unipotent, of order $2$, of order $3$.
\end{enumerate}
\end{namedthm}

\begin{namedcor}\label{corB}
\Proved{} The $\Sym^1$ square roots of Cooper's $s_{10}$ and $s_{18}$
\textup{(}$\delta_\infty=\tfrac14$\textup{)} and the infinite family
$(\alpha,-\tfrac\alpha6,-\tfrac\alpha3,-27,63,-30)$, integral exactly for
$\alpha\equiv18\pmod{36}$ \textup{(}$\rho=\tfrac76\equiv\tfrac16$\textup{)}, are
not Picard--Fuchs systems of any elliptic surface, in any gauge. The second
family has unbounded positive score and Ap\'ery limit
$-\tfrac1{2\sqrt3}\arctan\tfrac{6\sqrt3}\alpha$: it is Pad\'e approximation, not
geometry.
\end{namedcor}

The third tier is computational and stands in for the missing converse.

\begin{namedthm}[the scans]\label{thmC}
\VerifiedR{boxes as stated} An exhaustive integrality scan of
all $26$ Kodaira-admissible normalisation classes over $1\le A\le3000$,
$|B|\le150$, $|C|\le\min(4000,12000/M^2)$ --- $1\,270\,065$ hits, integrality
prime-tested to $n=30$ and re-verified exactly to $n=300$ --- returns $28$
primitive, Casoratian-non-degenerate rows in the rigid classes, of which exactly
ten have an elliptic-surface local system (Zagier's six, the two root rows of
Theorem~\ref{thmA}(2), and two new rows on $I_1I_7\,II\,II$ with monodromy
$\Gamma_0(7)$ and on $I_3\,III\,III\,III$); in the family classes the integral
rows form infinite families; and the only rows anywhere in the scan with positive
score and $k\ge2$ are Ap\'ery's $\zeta(2)$ row $(+0.40606)$ and Beukers' row
$(+0.13920)$. An independent modular scan of $3\,036\,854$ eta-quotient pairs
$(t,F)$ on $\Gamma_0(N)$, $N\le60$, produces exactly two classes with
$|\lambda_2|<1$: Beukers' and the level-$8$ $T$ row.
\end{namedthm}

\noindent
See Theorem~\ref{thm:herfclass} and Corollary~\ref{cor:H6} for the details.

Two further structural results are used throughout and stated in full below: the
cusp-move theorem (Theorem~\ref{thm:move}), which makes ``which fibre sits at
$\infty$'' into a group action and supplies Theorem~\ref{thmB}(iii); and the
host-group theorem (Theorem~\ref{thm:N2}), by which a three-term row forces its
host to have exactly four special points, hence covolume index $\le12$, so that
inside $\PSL_2(\Z)$ there are exactly $28$ candidate groups, $16$ congruence and
$12$ non-congruence --- with Beauville's six recovered as the four-cusp case.
Beukers' Theorem~3 row is the unique positive-score row whose form is
non-congruence.

\subsection*{The honest statement}
Theorem~\ref{thmA} classifies the rows whose operator is \emph{assumed} to be of
geometric origin. An unconditional classification --- ``every integral solution is
on the list'' --- needs the converse, and the converse is exactly Zagier's
conjecture: it is the rank-two, four-singular-point case of the Bombieri--Dwork
conjecture that a globally nilpotent operator comes from geometry. One direction
of the reduction is proved here: by Theorem~\ref{thm:N2} a modular
parametrisation of a three-term row forces four special points, hence one of $28$
genus-zero groups, hence (with Theorems~\ref{thmA} and \ref{thm:N3}) the list. So
\emph{Zagier's conjecture for this class is equivalent to the statement that the
projective monodromy group of an integral three-term row is commensurable with
$\PSL_2(\Z)$} (Proposition~\ref{prop:reduction}). That statement is not
accessible by rigidity --- rank two with four singular points is not rigid in
Katz's sense, the accessory parameter being the modulus --- nor by the naive
implication ``globally nilpotent $\Rightarrow$ hypergeometric pullback'', for
which Krammer's Lam\'e example \cite{Krammer1996,MovasatiReiter} is the standing
warning; the best partial result we know is Beukers' $p$-adic analysis of Dwork's
accessory parameter problem in exactly this configuration \cite{Beukers2002},
which does not close it. \S\ref{sec:BD} gives the evidence in place of a proof.
\Open{}

\subsection*{Status discipline}
Every statement carries its evidence class: \Proved{} (a proof or proof sketch is
given), \VerifiedR{$R$} (an exact finite computation over the stated range $R$ ---
not a proof, and never silently upgraded), \Numerical{$D$} (a floating-point or
$p$-adic measurement to $D$ digits), \Open{}. \S\ref{sec:repro} names the source
documents and scripts for every computation.

% ===========================================================================
\section{The exponents of a second-order row}\label{sec:exps}
% ===========================================================================

Let $\tht=t\,d/dt$. Substituting $y=\sum_{n\ge0}u_nt^n$ into \eqref{eq:R2gen}
gives the Picard--Fuchs operator
\begin{equation}\label{eq:Lgen}
L=\tht^2-t\,P(\tht)+t^2Q(\tht+1)
 =t^2R(t)\,\partial_t^2+t\,S(t)\,\partial_t+\bigl(-ct+(d{+}e{+}f)t^2\bigr),
\end{equation}
with
\[
R(t)=1-at+dt^2,\qquad S(t)=1-(a+b)t+(3d+e)t^2 .
\]
The characteristic roots of the recurrence are $\lambda_i=1/t_i$, the roots of
$\lambda^2-a\lambda+d$, where $t_1,t_2$ are the roots of $R$.

\begin{lemma}[exponent dictionary]\label{lem:exps}
\Proved{} Assume $d\neq0$ and $a^2\neq4d$, so that $L$ has exactly four singular
points $0,t_1,t_2,\infty$. Put $T(t):=S(t)-tR'(t)=1-bt+(d+e)t^2$. Then the local
exponents are
\[
(0,0)\ \text{at }t=0;\qquad
\Bigl(0,\ \rho_i=-\tfrac{T(t_i)}{t_iR'(t_i)}\Bigr)\ \text{at }t=t_i;\qquad
(1-r_1,\,1-r_2)\ \text{at }t=\infty,
\]
$r_1,r_2$ being the roots of $Q$; and $\rho_1+\rho_2=r_1+r_2=-e/d$
\textup{(}Fuchs\textup{)}. The point $t=0$ is a MUM point: the exponents are
equal and the second solution is logarithmic.
\end{lemma}

\begin{proof}
Expanding \eqref{eq:Lgen} gives the displayed $\partial_t$-form. At a simple zero
$t_i$ of $R$ the indicial polynomial is
$\rho\bigl[(\rho-1)t_iR'(t_i)+S(t_i)\bigr]$, whence the exponents $0$ and
$1-S(t_i)/(t_iR'(t_i))=-T(t_i)/(t_iR'(t_i))$. At $\infty$, $y\sim t^{\mu}$
forces $Q(\mu+1)=0$, so the exponents in $s=1/t$ are $-\mu=1-r_i$. The sum of the
six exponents is $2$, as it must be for a rank-two Fuchsian operator with four
singular points on $\PP^1$, and equals $\rho_1+\rho_2+2-(r_1+r_2)$.
\end{proof}

We write $\delta_\infty:=r_2-r_1$ for the exponent difference at $\infty$.

\begin{theorem}[normalisation classes]\label{thm:norm}
\Proved{} With the hypotheses of Lemma~\ref{lem:exps}, $\rho_1=\rho_2=\rho$ if
and only if
\begin{equation}\label{eq:normclass}
b=(1-\rho)a\qquad\text{and}\qquad e=-2\rho d .
\end{equation}
Writing $Q(n)=C(Mn-j_1)(Mn-j_2)$, so that $r_i=j_i/M$, this reads
\[
\rho=\frac{j_1+j_2}{2M},\qquad
\delta_\infty=\frac{j_2-j_1}{M},\qquad
P(n)=A\Bigl(n^2+\frac{2M-j_1-j_2}{2M}\,n\Bigr)+B ,
\]
with $A=a$, $B=c=u_1$ the accessory parameter and $D=CM^2=d$.
\end{theorem}

\begin{proof}
$\rho_1=\rho_2=\rho$ says $T(t_i)+\rho\,t_iR'(t_i)=0$ for $i=1,2$, i.e.\
$R\mid T+\rho tR'$. Both are quadratics with constant term $1$, so they are
equal; comparing coefficients gives $b+\rho a=a$ and $e+2\rho d=0$.
\end{proof}

Modulo the scaling $t\mapsto t/\mu$, $u_n\mapsto\mu^nu_n$ (which sends
$(a,b,c,d,e,f)\mapsto(\mu a,\mu b,\mu c,\mu^2d,\mu^2e,\mu^2f)$), a normalisation
class is fixed by the pair $(\rho,\delta_\infty)$; inside it the free parameters
are $A$, the leading coefficient $C$ of $Q$, and the accessory parameter
$B=u_1$. A rank-two Fuchsian equation with four singular points has exactly one
accessory parameter, and this is it.

Two special cases recover the two normalisations that occur in the literature.

\begin{itemize}[leftmargin=1.6em]
\item $\rho=0$ and $\delta_\infty=0$ give $b=a$, $e=f=0$: \emph{Zagier's shape}
$(an^2+an+c,\ dn^2)$. Geometrically $\rho=0$ says both movable fibres are
semistable and $\delta_\infty=0$ says the fibre at $\infty$ is too. So Zagier's
normalisation is not a choice: it is the statement that all four fibres are
semistable.
\item $\rho=\tfrac12$ and $\delta_\infty=0$ give $b=\tfrac a2$ and
$Q=C(2n-1)^2$: the \emph{root-row shape}, observed empirically for all nine
$\Sym^1$ square roots of the third-order sporadic families. Geometrically the two
movable fibres have order-two monodromy ($III$ or $III^*$) and the fibre at
$\infty$ has half-integral exponents ($I_m^*$).
\end{itemize}

\subsection{Kodaira admissibility}

\begin{theorem}\label{thm:kodaira}
\Proved{} Suppose the local system $\ker L$ is, up to a rational gauge factor and
a quadratic character twist, $R^1\pi_*\Q$ of an elliptic surface
$\pi:\mathcal E\to\PP^1_t$. Then
\[
\rho_1,\ \rho_2,\ \delta_\infty\ \in\ \{0,\tfrac12,\tfrac13,\tfrac23\}\pmod 1 .
\]
\end{theorem}

\begin{proof}
A rational gauge factor shifts exponents by integers, a quadratic twist by
half-integers, at each singular point; neither changes an exponent
\emph{difference} mod $1$. The exponent difference at a point is
$\tfrac1{2\pi i}\log$ of the ratio of the eigenvalues of the local monodromy, and
Kodaira's list \cite{Kodaira1963} gives the eigenvalue pairs $(1,1)$ for $I_n$,
$(-1,-1)$ for $I_n^*$, $(\pm i)$ for $III,III^*$, $(e^{\pm2\pi i/3})$ for
$IV,IV^*$ and $(e^{\pm\pi i/3})$ for $II,II^*$; the ratios are
$1,\,1,\,-1,\,e^{\mp4\pi i/3},\,e^{\mp2\pi i/3}$, i.e.\ differences
$0,0,\tfrac12,\tfrac13,\tfrac23$ mod $1$.
\end{proof}

The three admissible values are exactly the three kinds of point of a Fuchsian
group commensurable with $\PSL_2(\Z)$: a cusp, an elliptic point of order two,
an elliptic point of order three. The MUM point is always a cusp.

\begin{corollary}\label{cor:excluded}
\Proved{} The following are \emph{not} Picard--Fuchs systems of an elliptic
surface, in any gauge:
\begin{enumerate}[label=\textup{(\alph*)},leftmargin=2.4em]
\item the $\Sym^1$ square roots of Cooper's $s_{10}$ and $s_{18}$, whose trailing
coefficients are $-4(8n-3)(8n-5)$ and $12(8n-3)(8n-5)$, so that
$\delta_\infty=\tfrac14$ --- an order-eight local monodromy;
\item the one-parameter family
$(\alpha,-\tfrac\alpha6,-\tfrac\alpha3,-27,63,-30)$, integral exactly for
$\alpha\equiv18\pmod{36}$, which has $\rho=\tfrac76\equiv\tfrac16$.
\end{enumerate}
\end{corollary}

Item (b) is the sharpest available warning against reading a positive score as
evidence of a new object: the family is
$(n+1)^2u_{n+1}=(3n-2)\bigl[\tfrac\alpha6(2n+1)u_n+3(3n-5)u_{n-1}\bigr]$ with
$\lambda_1\lambda_2=-27$ and score $\approx\log(\alpha/27)-2\to\infty$, its member
$\alpha=234$ scoring $+0.15998$, above Beukers' row; and its Ap\'ery limit is
$-\tfrac1{2\sqrt3}\arctan\tfrac{6\sqrt3}\alpha$
\Numerical{193 digits, $\alpha=18,54,90,126,234,306,558,846$}, the $\arctan$
analogue of the classical Legendre--Pad\'e $\log$ family. Pad\'e approximation,
not geometry --- and Theorem~\ref{thm:kodaira} says so independently of the closed
form.

\begin{corollary}[Galois rigidity]\label{cor:galois}
\Proved{} If $a^2-4d$ is not a square then $t_1,t_2$ are conjugate irrationals,
hence $\rho_1=\rho_2$ and the two movable fibres have the same Kodaira class.
Equivalently: a configuration whose two movable fibres differ in class requires
$a^2-4d\in\Q^{\times2}$, so $\lambda_1,\lambda_2\in\Z$, so $|\lambda_2|\ge1$ and
the score $\log(1/|\lambda_2|)-k$ is $\le-k<0$.
\end{corollary}

\begin{proof}
Let $\sigma$ be the nontrivial automorphism of $\Q(t_1)$. The conditions
$T(t_i)+\rho_i t_iR'(t_i)=0$ have rational coefficients, so applying $\sigma$ to
the condition at $t_1$ gives $T(t_2)+\rho_1t_2R'(t_2)=0$; since
$t_2R'(t_2)\ne0$ this forces $\rho_2=\rho_1$. For the second clause,
$\lambda_{1,2}$ are the roots of the monic integer quadratic
$\lambda^2-a\lambda+d$, hence rational integers when the discriminant is a
square; the only integer of absolute value $<1$ is $0$, and $\lambda_2=0$ means
$d=0$, excluded.
\end{proof}

% ===========================================================================
\section{Herfurtner's list and the cross-ratio invariant}\label{sec:herf}
% ===========================================================================

Herfurtner \cite{Herfurtner1991} classifies the minimal elliptic surfaces over
$\PP^1$ with section, non-constant $\mathcal J$, and exactly four singular
fibres. Since $e(I_0^*)=6$, two fibres in $T^-=\{I_n^*,IV^*,III^*,II^*\}$ force
$\sum e\ge14>12$, so a \emph{rational} such surface has at most one fibre in
$T^-$; his Table~3 lists $50$ twist classes, and ``transfer of $*$''
($I_n\leftrightarrow I_n^*$, $II\leftrightarrow IV^*$, $III\leftrightarrow III^*$,
$IV\leftrightarrow II^*$) gives the full list of $56$ quadruples: $38$
\emph{rigid} (all fibres in $T^+$; the surface unique, $\mathcal J$ one of $38$
Belyi maps --- the same list as Vidunas--Filipuk's Table~4 \cite{VidunasFilipuk}
and Movasati--Reiter's Table~1 \cite{MovasatiReiter}) and $18$ in
\emph{one-parameter} families. One has $\deg\mathcal J=\sum n_i$ over the $I_n$
and $I_n^*$ fibres, and the six semistable configurations
$I_9I_1I_1I_1$, $I_8I_2I_1I_1$, $I_6I_3I_2I_1$, $I_5I_5I_1I_1$, $I_4I_4I_2I_2$,
$I_3I_3I_3I_3$ are Beauville's \cite{Beauville1982}, all rigid with
$\deg\mathcal J=12$.

A row imposes three conditions on a configuration: a cusp at $t=0$ (the MUM
point); the two movable points $t_1,t_2$ carry fibres of the \emph{same} Kodaira
class unless $a^2-4d$ is a square (Corollary~\ref{cor:galois}); and the fourth
fibre sits at $t=\infty$. Grouping the $38$ rigid configurations by the pair
(Kodaira class at $t_1,t_2$; class at $\infty$) gives nine normalisation
classes. Beauville's six occupy
(cusp,\,cusp), i.e.\ $(\rho;\delta_\infty)=(0;0)$; (cusp,\,ord\,$2$) holds
\#23--\#26 ($I_1I_1I_7III$, $I_1I_2I_6III$, $I_1I_3I_5III$, $I_2I_3I_4III$);
(cusp,\,ord\,$3$) holds \#19--\#22 and \#27--\#29; (ord\,$2$,\,cusp) --- the
root-row class --- holds \#38 $I_1I_5III\,III$, \#39 $I_2I_4III\,III$, \#40
$I_3I_3III\,III$; (ord\,$2$,\,ord\,$2$) holds only \#45 $I_3\,III\,III\,III$;
(ord\,$2$,\,ord\,$3$) holds \#44, \#48; (ord\,$3$,\,cusp) holds \#30
$I_1I_7II\,II$, \#31, \#32, \#36, \#37, \#43; (ord\,$3$,\,ord\,$2$) holds \#46,
\#49; and (ord\,$3$,\,ord\,$3$) holds \#47, \#50. \Proved{} (given Herfurtner's
table).

The one-parameter families enter as follows. An $I_0^*$ fibre has trivial image
in $\PSL_2$, so at such a point $L$ has an \emph{apparent} singularity (exponents
$(0,m)$ with $m\in\Z_{>0}$ and no logarithm): these are the classes with
$\rho\in\Z_{>0}$. An $I_m^*$ has half-integral exponents, hence in our
normalisation can only sit at $\infty$, and that is precisely what $Q=C(2n-1)^2$
does. So the root-row class $(\tfrac12;0)$ meets both the rigid block and, a
priori, families with an $I_m^*$ at $\infty$; but
$\sum e=I_a+3+3+(m+6)=12$ forces $a+m=0$, impossible, so in the root-row class
only \#38, \#39, \#40 survive.

\subsection{The cross-ratio invariant}

\begin{definition}\label{def:crossinv}
The row \eqref{eq:R2gen} has four singular points $0,t_1,t_2,\infty$, of which
$0$ is the MUM cusp and $\{t_1,t_2\}$ is an unordered pair. Its
\emph{cross-ratio invariant} is
\[
\mathcal I:=\frac{(1+z)^2}{z}=\frac{(\lambda_1+\lambda_2)^2}{\lambda_1\lambda_2}
=\frac{a^2}{d},\qquad z=\frac{t_1}{t_2}=\frac{\lambda_2}{\lambda_1},
\]
invariant under $z\mapsto1/z$, hence under both allowed relabellings
($t_1\leftrightarrow t_2$, and $0\leftrightarrow\infty$ when both are cusps).
\end{definition}

For a Herfurtner configuration with base points $(p_1,\dots,p_4)$ and an
assignment $(p_0,p_\infty\mid p',p'')$ the same quantity is
$\mathcal I=(1+z)^2/z$ with
$z=\frac{(p'-p_0)(p''-p_\infty)}{(p'-p_\infty)(p''-p_0)}$.

\begin{theorem}[finiteness]\label{thm:finite}
\Proved{} \textup{(}given Herfurtner's table\textup{)} A row \eqref{eq:R2gen}
with $a,d\in\Z$ is the Picard--Fuchs system of a \emph{rigid} Herfurtner
configuration only if $\mathcal I=a^2/d$ equals the cross-ratio invariant of that
configuration under an assignment compatible with Lemma~\ref{lem:exps}. Since
$\mathcal I\in\Q$, every configuration whose invariant is irrational or non-real
is excluded outright. Per class the admissible sets are finite and explicit; in
particular
\[
\text{all four fibres semistable:}\quad
\mathcal I\in\Bigl\{0,\ 3,\ \tfrac92,\ -\tfrac{49}8,\ \tfrac{100}9,\
\tfrac{289}{72},\ -121\Bigr\},
\]
\[
\text{root-row class }(\tfrac12,0):\quad
\mathcal I\in\Bigl\{\tfrac{484}{125}\ (\#38),\ 0\ (\#39),\ -12\ (\#40)\Bigr\},
\qquad
(\tfrac12;\tfrac12):\ \mathcal I=3\ (\#45),
\]
\[
(\tfrac13;0),(\tfrac23;0):\quad
\mathcal I\in\Bigl\{\tfrac{169}{49},\ 0,\ -50,\ \tfrac{1681}{400},\
-\tfrac{49}8\Bigr\}.
\]
\end{theorem}

\begin{proof}
$\mathcal I$ is a M\"obius invariant of the labelled quadruple, and the
identification of $\ker L$ with the surface's local system is by a M\"obius map
of $\PP^1_t$ matching the four singular points to the four base points, in a way
compatible with the Kodaira classes read off from Lemma~\ref{lem:exps}. For each
of the finitely many rigid configurations and each of the finitely many
compatible assignments one computes $\mathcal I$ from the base points; the
displayed lists are the outcome.
\end{proof}

The eighth semistable value is $65+\tfrac{682}{25}\sqrt5$, belonging to
$I_5I_5I_1I_1$ with the two $I_1$'s taken as the movable pair; it is irrational,
hence unreachable by an integer recurrence.

\subsection{Zagier's six, read off the table}

The six rows $\mathbf A(7,2,-8)$, $\mathbf B(9,3,27)$, $\mathbf C(10,3,9)$,
$\mathbf D(11,3,-1)$, $\mathbf E(12,4,32)$, $\mathbf F(17,6,72)$ have
$\mathcal I=a^2/c=-\tfrac{49}8,\,3,\,\tfrac{100}9,\,-121,\,\tfrac92,\,
\tfrac{289}{72}$, matching \#15 $I_1I_2I_3I_6$ (three times), \#13
$I_1I_1I_1I_9$, \#16 $I_1I_1I_5I_5$ and \#14 $I_1I_1I_2I_8$, all with MUM fibre
$I_1$ and $\deg\mathcal J=12$ (Table~\ref{tab:placements}). The matching is
Zagier's \cite[\S7]{Zagier2009}, obtained there by computing the $j$-invariant of
each Beauville family as a rational function of $t$ and inverting (the
calculations are in Verrill \cite{Verrill}); what the cross-ratio adds is that
the match can be made, and its completeness checked, without computing $j$ at
all. The invariant alone does not separate \#13 from \#18 (both $\mathcal I=3$)
nor \#14 from \#17 (both $\tfrac92$); the $\mathcal J$-map test of \S\ref{sec:J}
does, by returning the width $h$ of the MUM cusp: $h=1$ for $\mathbf B$ (so the
MUM fibre is $I_1$, hence \#13, since \#18 $=I_3I_3I_3I_3$ has all cusps of width
$3$) and $h=1$ for $\mathbf E$ (\#14, not \#17 whose widths are $2,2,4,4$).
\VerifiedR{cross-ratios exact; $\mathcal J$-fits exact}

The most interesting line is that \#15 occurs three times: $\mathbf A$,
$\mathbf C$, $\mathbf F$ are three rows on the \emph{same} surface with the
\emph{same} MUM fibre, differing only in which of $I_2,I_3,I_6$ is placed at
$\infty$. This is Zagier's table entry ``$IV\mapsto\mathbf A,\mathbf C,\mathbf
F$'' together with his remark
$t_{\mathbf C}=t_{\mathbf A}/(1+t_{\mathbf A})$; \S\ref{sec:cusp} makes it an
identity of operators with an explicit gauge, adds the second placement pair
$\{\mathbf E,\mathbf G\}$, proves the placement list complete, and --- this part
is not in \cite{Zagier2009} --- shows that the three rows on the one surface carry
\emph{two} different extension classes.

\subsection{Why ``six''}\label{sec:whysix}

Read the classification backwards. Fix all four fibres semistable. Then
Theorem~\ref{thm:norm} forces the shape $(an^2+an+b,\;cn^2)$;
Herfurtner (here $=$ Beauville) gives six surfaces, of which $V$ and $VI$ are
isogenous to $II$ and $I$ and give the same equations \cite[\S7]{Zagier2009};
Theorem~\ref{thm:finite} converts each configuration and placement into a value
of $a^2/c$; seven of the eight available values are rational, and each rational
value is realised by exactly one integral row. The eighth,
$65+\tfrac{682}{25}\sqrt5$, is irrational and therefore unreachable by an
integer recurrence. The value $\mathcal I=0$ needs $a=0$ and gives Zagier's
$\mathbf G$: $u_{2m}=\binom{2m}m^2$, $u_{2m+1}=0$, whose characteristic roots
$\pm4$ have equal modulus so that there is no Ap\'ery limit --- which is exactly
why Zagier's \emph{sporadic} list has six entries and the placement list has
seven. Hence:

\begin{quote}
\emph{The ``six'' of Beauville--Zagier is, on the recurrence side, the number of
rational cross-ratios in the four-cusp block of Herfurtner's table that carry a
non-degenerate integral accessory parameter.}
\end{quote}

The two Beauville configurations that do not appear as $I$--$VI$ labels in
Table~\ref{tab:placements} are $I_4I_4I_2I_2$ and $I_3I_3I_3I_3$; they are the
isogenous partners $V$ and $I$ of $II$ and $VI$, their only rational invariants
are $\mathcal I=0$ and $\mathcal I=3$, and those values are already taken by
$\mathbf G$ and $\mathbf B$.

% ===========================================================================
\section{The functional invariant and the $K3$ question}\label{sec:J}
% ===========================================================================

The cross-ratio test is necessary but not sufficient: it uses the positions of
the singular points and ignores the accessory parameter. A sufficient test comes
from the nome.

\begin{lemma}[the $\mathcal J$-map test]\label{lem:jtest}
\Proved{} Let $q=t\exp(g/y_0)$ be the canonical nome of $L$ at the MUM point,
where $y_1=y_0\log t+g$ is the second Frobenius solution. If $t$ is a Hauptmodul
of a genus-zero group $\Gamma$ with $t=c\,q_h+O(q_h^2)$ at a cusp of width $h$,
$q_h=e^{2\pi i\tau/h}$, then $q=c\,q_h$, hence $e^{2\pi i\tau}=(q/c)^h$.
Therefore $\Gamma$ is conjugate into $\PSL_2(\Z)$ --- equivalently $L$ is
projectively the Picard--Fuchs system of an elliptic surface over $\PP^1_t$ ---
if and only if for some $h\in\Z_{>0}$ and some constant $\gamma=c^{-h}$ the
composite $J\bigl(\gamma q(t)^h\bigr)$, with $J=q^{-1}+744+196884q+\cdots$ the
classical modular function, is a rational function of $t$; and then
$\deg\mathcal J=\deg J(t)$.
\end{lemma}

In practice one forms $W:=t^{h}J(\gamma q(t)^h)\in\Q[[t]]$ and solves
$V\cdot W=t^{h}U$ with $\deg U,\deg V\le d_{\max}$ by linear algebra over $\Q$,
checking the fit on more coefficients than were used to produce it. A
\emph{positive} verdict is a certificate; a \emph{negative} verdict is only ``no
fit within the stated window'', which is why the next proposition matters.

\begin{proposition}[the window is never the constraint]\label{prop:K1}
\Proved{} Let $\pi:\mathcal E\to\PP^1_t$ be a minimal elliptic surface with
section and non-constant functional invariant $\mathcal J$, $\mu:=\deg\mathcal J$,
whose singular fibres sit over the four singular points $0,t_1,t_2,\infty$ of the
row's operator. Among those four points let $c$ be the number carrying an
$I_n$ or $I_n^*$ fibre \textup{(}$n\ge1$; a cusp\textup{)}, $e_2$ the number
carrying $III$ or $III^*$, $e_3$ the number carrying $II,IV,IV^*,II^*$. Then
\[
\deg\mathcal J\ \le\ 6c+4e_3+3e_2-12\ \le\ 12 ,
\]
with equality on the right only in the all-cusp class. This holds for a surface
of any Kodaira dimension --- rational, $K3$, or higher.
\end{proposition}

\begin{proof}
By Kodaira's dictionary for the functional invariant (see Miranda
\cite{MirandaBook}, ch.~IV): $\mathcal J^{-1}(\infty)$ is exactly the set of
$I_n/I_n^*$ fibres, with ramification index $n$, so
$\#\mathcal J^{-1}(\infty)=c$ and $\sum_{\mathcal J(P)=\infty}(e_P-1)=\mu-c$;
over $j=0$ the ramification index is divisible by $3$ except at fibres of type
$II,IV,IV^*,II^*$, so at most $e_3$ preimages have $e_P<3$ and
$\sum_{\mathcal J(P)=0}(e_P-1)\ge\tfrac23(\mu-e_3)$; over $j=1728$ the index is
even except at $III,III^*$, giving
$\sum_{\mathcal J(P)=1728}(e_P-1)\ge\tfrac12(\mu-e_2)$. Riemann--Hurwitz for
$\mathcal J:\PP^1\to\PP^1$ gives
\[
2\mu-2=\sum_P(e_P-1)\ \ge\ (\mu-c)+\tfrac23(\mu-e_3)+\tfrac12(\mu-e_2)
=\tfrac{13}6\mu-c-\tfrac23e_3-\tfrac12e_2 ,
\]
i.e.\ $\tfrac16\mu\le c+\tfrac23e_3+\tfrac12e_2-2$, which is the bound. Any
further ramification only strengthens it. Maximising $6c+4e_3+3e_2-12$ subject to
$c+e_2+e_3\le4$ gives $12$, attained only at $c=4$. Fibres of type $I_0^*$ have
trivial image in $\PSL_2$, are invisible to $L$, and change nothing.
\end{proof}

This is the genus formula $g=1+\tfrac\mu{12}-\tfrac{e_2}4-\tfrac{e_3}3-\tfrac c2$
for a genus-zero Fuchsian group of index $\mu$ and signature
$(0;2^{e_2},3^{e_3};c)$, written so that it survives extra ramification, i.e.\ so
that it does not assume $t$ is a Hauptmodul. Class by class the bound reads
$\le12$ for $(0;0)$, $\le6$ for $(\tfrac12;0)$, $\le8$ for $(\tfrac13;0)$,
$\le4$ for $(\tfrac12;\tfrac13)$, $\le3$ for $(\tfrac12;\tfrac12)$; and for the
three classes whose movable points carry an apparent singularity
($\rho\in\Z_{>0}$) the bound is $\le0$, i.e.\ those classes contain no
elliptic-surface row at all, independently of any search.

Riemann--Hurwitz is not the whole story: it bounds $\deg\mathcal J$ but says
nothing about $\chi$, and a quadratic twist changes $\chi$ without changing
$\mathcal J$. The missing constraint is a parity.

\begin{lemma}[twist parity]\label{lem:twistparity}
\Proved{} In the situation of Proposition~\ref{prop:K1}, let $e_{\min}(P)$ be the
Euler number of the \emph{unstarred} member of the twist pair at $P$: $n$ at an
$I_n/I_n^*$, $3$ at a $III/III^*$, $2$ at a $II/IV^*$, $4$ at a $IV/II^*$. Then
\[
12\ \mid\ \sum_{P}e_{\min}(P)\ =\ \deg\mathcal J+3e_2+\!\!\sum_{\text{ord-}3\ P}\!\!e_{\min}(P),
\]
the sum running over the four singular points of the row.
\end{lemma}

\begin{proof}
For a relatively minimal elliptic surface with section over $\PP^1$,
$\sum e=12\chi$. Every singular fibre of $\mathcal E$ lies either over one of the
four points or is an $I_0^*$ (trivial image in $\PSL_2$, invisible to $L$; any
$I_n^*$ with $n\ge1$ elsewhere would be a fifth singular point of $L$). Fix the
reference twist in which every fibre is unstarred; the actual surface is its
quadratic twist by a character $\varepsilon_S$ with branch set $S$, and at each point of
$S$ the fibre flips to its starred partner, raising $e$ by exactly $6$ --- for
$I_n\!\to\!I_n^*$, $III\!\to\!III^*$, $II\!\to\!IV^*$, $IV\!\to\!II^*$ and
$I_0\!\to\!I_0^*$ alike. Hence $\sum e=\sum_Pe_{\min}(P)+6\,|S|$. A quadratic
character on $\PP^1$ has an even number of branch points, so
$\sum e\equiv\sum_Pe_{\min}(P)\pmod{12}$, and $12\mid\sum e$. Finally
$\sum_Pe_{\min}(P)=\deg\mathcal J+3e_2+\sum_{\text{ord-}3}e_{\min}$, since the
cusps contribute their $n_i$, whose sum is $\deg\mathcal J$.
\end{proof}

Proposition~\ref{prop:K1} and Lemma~\ref{lem:twistparity} together pin
$\deg\mathcal J$ exactly in every class, and pin the surface up to quadratic
twist. Class by class:

\smallskip
{\small
\begin{tabular}{@{}llllll@{}}
\hline
class $(\rho;\delta_\infty)$ & $(c,e_2,e_3)$ & RH bound & parity & forced $\deg\mathcal J$ & untwisted configuration\\
\hline
$(0;0)$ & $(4,0,0)$ & $\le12$ & $12\mid D$ & $12$ & Beauville's six\\
$(\tfrac12;0)$ & $(2,2,0)$ & $\le6$ & $12\mid D{+}6$ & $6$ & \#38, \#39, \#40\\
$(\tfrac13;0)$, $(\tfrac23;0)$ & $(2,0,2)$ & $\le8$ & $12\mid D{+}s$, $s\in\{4,6,8\}$ & $8,6,4$ & \#30--\#32; \#36, \#37; \#43\\
$(\tfrac12;\tfrac12)$ & $(1,3,0)$ & $\le3$ & $12\mid D{+}9$ & $3$ & \#45\\
$(\tfrac12;\tfrac13)$, $(\tfrac12;\tfrac23)$ & $(1,2,1)$ & $\le4$ & $12\mid D{+}6{+}e_{\min}$ & $4$ or $2$ & \#48; \#44\\
\hline
\end{tabular}
}
\smallskip

\noindent
In every case the untwisted configuration has $\sum e=12$, i.e.\ $\chi=1$: it is
a \emph{rational} elliptic surface. In the third row $s=e_{\min}(t_1)+e_{\min}(t_2)$
is $4$, $6$ or $8$ according as $t_1,t_2$ carry $II,II$; $II,IV$; or $IV,IV$, and
by Corollary~\ref{cor:galois} the mixed case $s=6$ needs $a^2-4d$ a square. The
point of the lemma is that the parity, not Riemann--Hurwitz, is what excludes the
smaller degrees: in the all-cusp class RH alone permits $\deg\mathcal J=6$ (with
one extra simple ramification point, i.e.\ a one-parameter family carrying a
single $I^*$ fibre), and it is $12\mid\deg\mathcal J$ that rules it out.

\begin{remark}[the $K3$ recheck]
An elliptic $K3$ with four singular fibres has $\sum e=24$, so one might expect
$\deg\mathcal J$ up to $24$. Proposition~\ref{prop:K1} says otherwise: the
difference between $\sum e=12$ and $24$ is absorbed by fibres contributing to $e$
but not to $\deg\mathcal J$ ($II,III,IV,IV^*,III^*,II^*$ contribute
$2,3,4,8,9,10$ to $\sum e$ and $0$ to $\deg\mathcal J$), not by a larger
$\mathcal J$; equivalently, by Lemma~\ref{lem:twistparity} a $K3$ realisation is
a quadratic twist of a rational one and shares its $\mathcal J$. The computation agrees: re-run with $\deg\mathcal J\le24$,
$h\le24$, $140$ series terms and $992$ values of $\gamma$ per $h$ (against $12$,
$12$, $56$, $96$) on all $18$ negative rows, on three Kodaira-inadmissible
controls and on the six known positives, \emph{every negative stays negative and
every positive returns the same $(h,\deg\mathcal J)$}, with $50$ coefficients
fitted and $65$ further identities verified.
\VerifiedR{$\deg\le24$, $h\le24$, $140$ terms, $23\,808$ $(h,\gamma)$ pairs per
row, $27$ rows}
\end{remark}

\begin{theorem}[which root rows are elliptic surfaces]\label{thm:whichroots}
\VerifiedR{exact, series order $56$ resp.\ $140$, degrees $\le12$ resp.\ $\le24$}
Of the nine $\Sym^1$ square roots of the third-order sporadic families, exactly
two are Picard--Fuchs systems of an elliptic surface over their own $\PP^1_t$:
\[
\sqrt{\mathrm{AZ}(11,5,125)}\ \leftrightarrow\ I_1I_5\,III\,III\ (\#38,\ \deg\mathcal J=6,\ h=1),
\]
\[
\sqrt{\mathrm{AZ}(9,3,-27)}\ \leftrightarrow\ I_3I_3\,III\,III\ (\#40,\ \deg\mathcal J=6,\ h=3),
\]
in agreement with $\mathcal I=\tfrac{484}{125}$ and $\mathcal I=-12$. The
remaining seven --- Ap\'ery's own square root, i.e.\ Beukers' Theorem~3 row
($\mathcal I=1156$), the $T$ row ($36$), Domb's ($\tfrac{25}4$),
$\mathrm{AZ}(7,3,81)$'s ($\tfrac{196}{81}$) and Cooper's $s_7$'s
($-\tfrac{676}{27}$), together with $\sqrt{s_{10}},\sqrt{s_{18}}$ which are
excluded already by Corollary~\ref{cor:excluded} --- fail the test: their
projective monodromy is an Atkin--Lehner quotient $\Gamma_0(N)+W$, commensurable
with but not conjugate into $\PSL_2(\Z)$.
\end{theorem}

The mechanism is transparent on Ap\'ery's curve, where
$t=(\eta_1\eta_6/\eta_2\eta_3)^{12}$ has degree two on $X_0(6)$ and is
$W_6$-invariant, so $\PP^1_t=X_0(6)/W_6$: the universal elliptic curve does not
descend along that quotient, only its symmetric square does, and $j$ is not
$W_6$-invariant, so no $\mathcal J$-map of finite degree exists. That is why
Ap\'ery's \emph{order-three} row lives happily on $\PP^1_t$ while its square root
does not. The two fixed points of $W_6$ become the two order-two points of the
quotient, i.e.\ the two exponents $\rho=\tfrac12$ that the root-row shape
records; one is the Fricke fold point $i/\sqrt6$, with
$t(i/\sqrt6)=(\sqrt2-1)^4$ exactly. Same story for $T$ ($W_8$ on $X_0(8)$), Domb
and $\mathrm{AZ}(7,3,81)$.

\begin{remark}
The $\mathcal J$-test is not implied by the cross-ratio test. Row \#14 of
Table~\ref{tab:rows} has $\mathcal I=-50$, the invariant of Herfurtner \#32
$I_4I_4\,II\,II$, yet the $\mathcal J$-test says no: the cross-ratio ignores the
accessory parameter. Making the test sufficient would require computing the
accessory parameter of each Herfurtner Weierstrass model from its $G_2,G_3$ and
comparing --- mechanical, and \Open{} here.
\end{remark}

% ===========================================================================
\section{The two-tier classification, and the scan}\label{sec:scan}
% ===========================================================================

\subsection{Proof of Theorem~\ref{thmA}}\label{sec:thmA}

\emph{Part \textup{(1)}.} The Zagier shape has $b=a$ and $e=f=0$, so by
Lemma~\ref{lem:exps} $\rho_1=\rho_2=0$ and $r_1=r_2=0$, i.e.\ all four exponent
differences vanish. By Theorem~\ref{thm:kodaira} the projective local monodromy
at each of the four points is unipotent, so each singular fibre is of type $I_n$
or $I_n^*$; this is (a). For (b), Proposition~\ref{prop:K1} with $(c,e_2,e_3)=(4,0,0)$ gives
$\deg\mathcal J\le12$, and Lemma~\ref{lem:twistparity} with $e_2=e_3=0$ gives
$12\mid\deg\mathcal J$; each of the four points carries a nontrivial unipotent
projective monodromy, so $n_i\ge1$ and $\deg\mathcal J=\sum n_i\ge4$. Hence
$\deg\mathcal J=12$. Equality in Riemann--Hurwitz then forces $\mathcal J$ to be
unramified outside $\{0,1728,\infty\}$ with the minimal indices there --- a Belyi
map --- and the untwisted member of the twist class has $\sum e=\sum n_i=12$,
i.e.\ $\chi=1$: a rational semistable elliptic surface with four singular fibres,
one of Beauville's six \cite{Beauville1982}. The surface itself is that one up to
quadratic twist, so $\chi=1+|S|/2$ may be $2$ (a $K3$) or more; but the twist
changes neither $\mathcal J$, nor the local system, nor the row, which is why no
hypothesis on the Kodaira dimension is needed and why the $K3$ range produces
nothing new. Riemann--Hurwitz alone would leave $\deg\mathcal J=6$ open, with one
extra simple ramification point and a single $I^*$ fibre --- exactly Herfurtner's
one-parameter families; it is the parity $12\mid\deg\mathcal J$ that excludes it.
For (c), the MUM point is a cusp and $t$ pulls back the canonical nome
(Lemma~\ref{lem:jtest}); normalising $t=q_h+O(q_h^2)$ makes $t$ a Hauptmodul with
integral $q$-expansion, and $u_n=[t^n]F$ for the weight-one form $F$ attached to
the fibration. For (d), $\mathcal I=a^2/c$ is a M\"obius invariant of the
labelled quadruple, so it must be the invariant of the configuration under the
chosen placement; Theorem~\ref{thm:finite} lists the seven rational values
available in the four-cusp block. Given a configuration and a placement --- which
cusp at $t=0$, which at $t=\infty$ --- the rigid surface, and hence its
Picard--Fuchs operator, is determined, so the accessory parameter $b$ is
determined; the enumeration of the resulting placements is
Table~\ref{tab:placements}.
\VerifiedR{cross-ratios exact; the $\mathcal J$-fits of \S\ref{sec:J} exact,
$56$-term series, degrees $\le12$, re-checked at $\deg\le24$} \qed

\begin{table}[t]
\small
\begin{tabular}{@{}llllll@{}}
\hline
row & Zagier \# & $(a,b,c)$ & configuration & fibre at $t=\infty$ & $\mathcal I=a^2/c$\\
\hline
$\mathbf A$ & \#5 & $(7,2,-8)$ & $I_6I_3I_2I_1$ ($IV$) & $I_2$ & $-49/8$\\
$\mathbf C$ & \#8 & $(10,3,9)$ & $I_6I_3I_2I_1$ ($IV$) & $I_3$ & $100/9$\\
$\mathbf F$ & \#13 & $(17,6,72)$ & $I_6I_3I_2I_1$ ($IV$) & $I_6$ & $289/72$\\
$\mathbf B$ & \#7 & $(9,3,27)$ & $I_1I_1I_1I_9$ ($I,VI$) & $I_1$ & $3$\\
$\mathbf D$ & \#9 & $(11,3,-1)$ & $I_1I_1I_5I_5$ ($III$) & $I_5$ & $-121$\\
$\mathbf E$ & \#10 & $(12,4,32)$ & $I_1I_1I_2I_8$ ($II,V$) & $I_2$ & $9/2$\\
$\mathbf G$ & \#2 & $(0,0,-16)$ & $I_1I_1I_2I_8$ ($II,V$) & $I_8$ & $0$\\
\hline
\end{tabular}
\medskip
\caption{The complete list of cusp placements in the Zagier shape
\textup{(}Theorem~\ref{thmA}(1)(d)\textup{)}, in Zagier's own labelling
$\mathbf A$--$\mathbf G$ and with his index numbers; the Roman numerals are
Beauville's family labels as in \cite[\S7]{Zagier2009}. The cusp-move orbits are
$\{\mathbf A,\mathbf C,\mathbf F\}$, $\{\mathbf B\}$, $\{\mathbf D\}$,
$\{\mathbf E,\mathbf G\}$ \textup{(}\S\ref{sec:cusp}\textup{)}. $\mathbf G$ is
$u_{2m}=\binom{2m}m^2$, $u_{2m+1}=0$: integral but with
$|\lambda_1|=|\lambda_2|=4$, hence no Ap\'ery limit, which is why Zagier's
\emph{sporadic} list stops at six. Zagier's remaining modular solutions \#11
$(16,4,0)$, \#14 $(27,6,0)$ and \#20 $(64,12,0)$ have $c=0$ and only three
singular points, so they fall outside the hypothesis and belong to
Schmickler-Hirzebruch's three-fibre classification
\cite{SchmicklerHirzebruch}.}\label{tab:placements}
\end{table}

\emph{Part \textup{(2)}.} The root-row shape has $b=\tfrac a2$ and
$Q=C(2n-1)^2$, so $\rho_1=\rho_2=\tfrac12$ and $r_1=r_2=\tfrac12$: the movable
points have exponent difference $\tfrac12$ and $\infty$ has exponents
$(\tfrac12,\tfrac12)$. By Theorem~\ref{thm:kodaira} the movable fibres are $III$
or $III^*$ and the fibre at $\infty$ is a cusp $I_m$ or $I_m^*$, with $m\ge1$:
$m=0$ would make $\infty$ an $I_0^*$, invisible to $L$, and
Proposition~\ref{prop:K1} with $(c,e_2,e_3)=(1,2,0)$ would then give
$\deg\mathcal J\le0$. Likewise $n\ge1$ at the MUM point. Now
Proposition~\ref{prop:K1} with $(c,e_2,e_3)=(2,2,0)$ gives
$\deg\mathcal J=n+m\le6\cdot2+3\cdot2-12=6$, while Lemma~\ref{lem:twistparity}
gives $12\mid n+m+6$, i.e.\ $\deg\mathcal J\equiv6\pmod{12}$. With
$2\le\deg\mathcal J\le6$ this forces $\deg\mathcal J=6$, and the untwisted member
of the twist class has fibres $I_n,I_m,III,III$ with $n+m=6$ and
$\sum e=6+3+3=12$: a rational elliptic surface, the three possibilities
$(n,m)=(1,5),(2,4),(3,3)$ being Herfurtner's \#38, \#39, \#40 with cross-ratios
$\tfrac{484}{125},0,-12$.

It is worth being explicit about what Riemann--Hurwitz does \emph{not} do here.
Euler numbers alone leave $\chi=2$ open. Writing $\sigma$ for the number of stars
among $\{0,\infty\}$: with $III$ at both movable points
$\sum e=n+m+6+6\sigma=12\chi$, and $\deg\mathcal J=n+m\le6$ still admits
$(\sigma,\chi)=(2,2)$; with $III^*$ at both, $n+m=12\chi-24-6\sigma$ leaves
nothing new; and in the \emph{mixed} case $III$ at $t_1$, $III^*$ at $t_2$ ---
indistinguishable at the level of exponents, both having difference $\tfrac12$ ---
one gets $n+m=12\chi-18-6\sigma$, so $(\sigma,\chi)=(1,2)$ with $n+m=6$ is
permitted: an elliptic $K3$ with fibres $III$, $III^*$, $I_n$, $I_m^*$, $n+m=6$,
$\sum e=3+9+n+(m+6)=24$ and $\deg\mathcal J=6$, comfortably inside the bound.
\emph{Riemann--Hurwitz does not exclude it.} What removes it as a source of new
rows is that it is the quadratic twist with branch set $\{t_2,\infty\}$ of the
corresponding \#38/\#39/\#40 configuration --- twisting flips $III^*\!\to\!III$ at
$t_2$ and $I_m^*\!\to\!I_m$ at $\infty$, lowering $\sum e$ by $12$ to $\chi=1$.
The two surfaces have the same $\mathcal J$, the same projective monodromy, and
local systems differing by the quadratic twist already allowed in the hypothesis;
so they carry the same row. Lemma~\ref{lem:twistparity} is the systematic form of
this observation. It is also confirmed computationally: the $\mathcal J$-map test
run at $\deg\mathcal J\le24$, $h\le24$ finds no fit for any root row other than
the two on \#38 and \#40 \VerifiedR{$27$ rows, $23\,808$ $(h,\gamma)$ pairs
each}, so no $K3$ realisation occurs among the known root rows either.

The realisations, and the negative verdict for the other seven $\Sym^1$ square
roots, are Theorem~\ref{thm:whichroots}; the value $\mathcal I=0$ needs $A=0$,
and the scan below finds no non-degenerate integral row there. \qed

\subsection{The scan}

Theorem~\ref{thm:norm} turns the search for integral rows into $26$ independent
three-parameter searches, one for each Kodaira-admissible pair $(\rho,\delta_\infty)$ with $r_i=j_i/M\in[-1,2]$, namely
$(M;j_1,j_2)$ running over
$(1;0,0)$, $(2;1,1)$, $(3;1,2)$, $(4;1,3)$, $(6;1,5)$, $(3;1,1)$, $(3;2,2)$,
$(1;0,1)$, $(1;1,1)$, $(2;1,3)$, $(2;-1,1)$, $(1;-1,1)$, $(3;0,2)$, $(3;1,3)$,
$(3;-1,1)$, $(4;-1,1)$, $(6;1,3)$, $(6;3,5)$, $(6;-1,1)$, $(12;1,7)$,
$(12;5,11)$, $(2;3,3)$, $(1;1,2)$, $(3;2,4)$, $(4;3,5)$, $(6;5,7)$
in the notation $Q(n)=C(Mn-j_1)(Mn-j_2)$ of Theorem~\ref{thm:norm}. Inside a
class the free integers are $A$, the accessory parameter $B=u_1$, and $C$. The
scan sweeps $1\le A\le3000$ (with the divisibility $2M\mid A(2M-j_1-j_2)$),
$1\le|B|\le150$, $1\le|C|\le\min(4000,12000/M^2)$, plus a separate $A=0$ pass,
testing integrality division-free on $U_n=u_n(n!)^2$,
\[
U_{n+1}=P(n)U_n-n^2Q(n)U_{n-1},\qquad u_n\in\Z\iff v_p(U_n)\ge2v_p(n!),
\]
prime by prime modulo $p^K$ for every $p\le30$. Survivors are re-verified with
exact rational arithmetic to $n=300$; the sharp denominator exponent $k$ of the
companion ($b_0=0$, $b_1=1$) is measured to $n=60$. The sweep is
$1\,270\,065$ hits in about $75$ minutes on twelve cores.

\subsection{Two filters}
The Casoratian satisfies
$u_nb_{n+1}-u_{n+1}b_n=\pm\prod_{m\le n}Q(m)/((m+1)!)^2$, so it vanishes
identically as soon as $Q(n_0)=0$ for an integer $n_0\ge1$ (i.e.\ $M\mid j_i$
with $j_i/M\ge1$); the companion is then a rational multiple of the row, the
Ap\'ery limit is rational, and the extension class is split. And $a^2=4d$ gives
$\lambda_1=\lambda_2$, hence no second growth rate. Both are applied before any
row is called interesting. The Casoratian filter is not cosmetic: in one early
run it reduced $25\,390$ raw hits to $248$, and without it the entire top of the
raw ranking consists of degeneracies.

\subsection{Rigid classes and family classes}
The scan output splits sharply, and the split is Herfurtner's own. Classes
attached to rigid configurations yield a handful of rows each ($10$--$140$ hits);
two classes yield tens of thousands, namely $(2;-1,1)$, i.e.\
$(\rho;\delta_\infty)=(0;1)$, with $28\,143$ hits, and $(2;1,3)$, i.e.\ $(1;1)$,
with $34\,430$ --- precisely the classes in which the fibre at $\infty$ has
half-integral exponents (an $I_m^*$), the second additionally having an apparent
singularity ($I_0^*$) at $t_1,t_2$. Those are exactly the Kodaira types that put
a configuration into Herfurtner's one-parameter block. For $(2;-1,1)$ the family
is explicit: $B=\tfrac A4$, $P(n)=\tfrac A4(2n+1)^2$, $Q(n)=C(2n-1)(2n+1)$,
integral for $A\in4\Z$ and $C$ in an arithmetic progression --- one parameter after
$t\mapsto t/\mu$, matching the one modulus of Herfurtner's $I_1I_1I_1I_3^*$ and
$I_1I_1I_2I_2^*$ families; its $C=4$ sub-family is the classical Legendre--Pad\'e
row with Ap\'ery limit $-\tfrac1{12}\log\frac{x+1}{x-1}$ at $A=8x$, which has
$k=1$, score $\to\infty$, and proves nothing new. \emph{This is the structural
dichotomy of the classification:} a rigid configuration supports at most finitely
many integral rows (Theorem~\ref{thm:finite} pins $a^2/d$ and the surface then
determines the accessory parameter), a one-parameter family supports an infinite
family of them, and Zagier's finiteness statement is a statement about the rigid
block.

\subsection{The rows}

\begin{table}[t]
\scriptsize
\begin{tabular}{@{}rlllrrrrl@{}}
\hline
\# & class & $P(n)$ & $Q(n)$ & $\lambda_1$ & $\lambda_2$ & $k$ & score & $\mathcal J$ / name\\
\hline
1 & $(1;0,0)$ & $9n^2{+}9n{+}3$ & $27n^2$ & $5.196$ & cplx & 2 & --- & yes $(1,12)$, Zagier $\mathbf B$\\
2 & $(2;1,1)$ & $56n^2{+}28n{+}6$ & $324(2n{-}1)^2$ & $36$ & cplx & 2 & --- & no, $\sqrt{\mathrm{AZ}(7,3,81)}$\\
3 & $(2;1,1)$ & $88n^2{+}44n{+}10$ & $500(2n{-}1)^2$ & $44.72$ & cplx & 2 & --- & \textbf{yes} $(1,6)$, $\sqrt{\mathrm{AZ}(11,5,125)}$\\
4 & $(3;1,1)$ & $117n^2{+}78n{+}21$ & $441(3n{-}1)^2$ & $63$ & cplx & 2 & --- & \textbf{yes} $(1,8)$, \textbf{NEW} $I_1I_7II\,II$\\
5 & $(3;1,2)$ & $40n^2{+}20n{+}4$ & $48(3n{-}1)(3n{-}2)$ & $20.78$ & cplx & 2 & --- & no, new\\
6 & $(4;1,3)$ & $72n^2{+}36n{+}6$ & $108(4n{-}1)(4n{-}3)$ & $41.57$ & cplx & 2 & --- & \textbf{yes} $(3,3)$, \textbf{NEW} $I_3III\,III\,III$\\
7 & $(1;0,0)$ & $11n^2{+}11n{+}3$ & $-n^2$ & $11.090$ & $-0.0902$ & 2 & $\mathbf{+0.406}$ & yes $(1,12)$, Zagier $\mathbf D$\\
8 & $(2;1,1)$ & $136n^2{+}68n{+}10$ & $4(2n{-}1)^2$ & $135.88$ & $0.1177$ & 2 & $\mathbf{+0.139}$ & no, Beukers $=\sqrt{\text{Ap\'ery}}$\\
9 & $(2;1,1)$ & $24n^2{+}12n{+}2$ & $4(2n{-}1)^2$ & $23.314$ & $0.6863$ & 2 & $-1.624$ & no, $\sqrt T$\\
10 & $(4;1,3)$ & $88n^2{+}44n{+}6$ & $-4(4n{-}1)(4n{-}3)$ & $88.721$ & $-0.7214$ & 2 & $-1.673$ & no, level-5 Fricke\\
11 & $(1;0,0)$ & $7n^2{+}7n{+}2$ & $-8n^2$ & $8$ & $-1$ & 2 & $-2.000$ & yes $(1,12)$, Zagier $\mathbf A$\\
12 & $(1;0,0)$ & $10n^2{+}10n{+}3$ & $9n^2$ & $9$ & $1$ & 2 & $-2.000$ & yes $(1,12)$, Zagier $\mathbf C$\\
13 & $(3;1,2)$ & $26n^2{+}13n{+}2$ & $-3(3n{-}1)(3n{-}2)$ & $27$ & $-1$ & 2 & $-2.000$ & no, $\sqrt{s_7}$\\
14 & $(3;1,1)$ & $180n^2{+}120n{+}24$ & $-72(3n{-}1)^2$ & $183.53$ & $-3.531$ & 2 & $-3.262$ & no ($\mathcal I=-50$), new\\
15 & $(1;0,0)$ & $12n^2{+}12n{+}4$ & $32n^2$ & $8$ & $4$ & 2 & $-3.386$ & yes $(1,12)$, Zagier $\mathbf E$\\
16 & $(2;1,1)$ & $20n^2{+}10n{+}2$ & $16(2n{-}1)^2$ & $16$ & $4$ & 2 & $-3.386$ & no, $\sqrt{\text{Domb}}$\\
17 & $(4;1,3)$ & $28n^2{+}14n{+}2$ & $-8(4n{-}1)(4n{-}3)$ & $32$ & $-4$ & 2 & $-3.386$ & no, new\\
18 & $(2;1,1)$ & $72n^2{+}36n{+}6$ & $-108(2n{-}1)^2$ & $77.569$ & $-5.569$ & 2 & $-3.717$ & \textbf{yes} $(3,6)$, $\sqrt{\mathrm{AZ}(9,3,-27)}$\\
19 & $(1;0,0)$ & $17n^2{+}17n{+}6$ & $72n^2$ & $9$ & $8$ & 2 & $-4.079$ & yes $(1,12)$, Zagier $\mathbf F$\\
20 & $(4;1,3)$ & $80n^2{+}40n{+}6$ & $36(4n{-}1)(4n{-}3)$ & $72$ & $8$ & 2 & $-4.079$ & no, new\\
21--22 & $(2;1,3)$ & $\pm4$ & $-64(2n{-}1)(2n{-}3)$ & $16$ & $-16$ & 2 & $-4.773$ & no, $A=0$\\
23 & $(4;1,3)$ & $48n^2{+}24n{+}4$ & $32(4n{-}1)(4n{-}3)$ & $32$ & $16$ & 2 & $-4.773$ & no, new\\
24--25 & $(3;2,4)$ & $\pm3$ & $-81(3n{-}2)(3n{-}4)$ & $27$ & $-27$ & 2 & $-5.296$ & no, $A=0$\\
26 & $(4;1,3)$ & $68n^2{+}34n{+}6$ & $72(4n{-}1)(4n{-}3)$ & $36$ & $32$ & 2 & $-5.466$ & no, new\\
27--28 & $(4;3,5)$ & $\pm4$ & $-256(4n{-}3)(4n{-}5)$ & $64$ & $-64$ & 2 & $-6.159$ & no, $A=0$\\
\hline
\end{tabular}
\medskip
\caption{All primitive, Casoratian-non-degenerate, non-double-root integral rows
produced by the scan in the classes attached to rigid configurations, sorted by
score. ``cplx'' means complex-conjugate characteristic roots, hence no
archimedean Ap\'ery limit. \VerifiedR{integrality exact to $n=300$; $k$ sharp to
$n=60$; $\mathcal J$-fits exact}}\label{tab:rows}
\end{table}

Table~\ref{tab:rows} is the answer in the rigid block. Four observations.

\begin{enumerate}[label=\textup{(\arabic*)},leftmargin=2.6em]
\item \emph{Every second-order row of the known census appears}, with its known
invariants: Zagier's six, the six Almkvist--Zudilin square roots, Cooper's
$\sqrt{s_7}$, the level-$5$ Fricke row. And $k=2$ for every row in the table
without exception.
\item \emph{Two rows are new and are elliptic surfaces} (rows 4 and 6);
see \S\ref{sec:newrows}.
\item Six further rows are new but fail the $\mathcal J$-test (rows 5, 14, 17,
20, 23, 26); like Beukers' row they live on Atkin--Lehner quotients.
\item \emph{No new row has a positive score.}
\end{enumerate}

\subsection{The two new elliptic-surface rows}\label{sec:newrows}

Both are invisible to Zagier's normalisation ($\rho\neq0$) \emph{and} to the
root-row normalisation, which is why no earlier scan found them. They are found
only by running the exponent classification first and scanning each
Kodaira-admissible class in turn.

\begin{proposition}\label{prop:newrows}
\VerifiedR{integrality $n\le300$ exact; $\mathcal J$-fit exact; uniformisation to
$q^{90}$ resp.\ $w^{45}$}
\begin{enumerate}[label=\textup{(\alph*)},leftmargin=2.4em]
\item $(n+1)^2u_{n+1}=(117n^2+78n+21)u_n-441(3n-1)^2u_{n-1}$, $u_0=1$,
\[u_n=1,\,21,\,693,\,23940,\,734643,\,13697019,\dots\in\Z ,\]
class $(\rho;\delta_\infty)=(\tfrac13;0)$, $\mathcal I=\tfrac{169}{49}$,
$\deg\mathcal J=8$, $h=1$: the surface $I_1I_7\,II\,II$ \textup{(}Herfurtner
\#30\textup{)}, projective monodromy $\Gamma_0(7)$. Uniformisation:
\[
t=-\tfrac19\Bigl(\frac{\eta(7\tau)}{\eta(\tau)}\Bigr)^4,\qquad
A(t)\,\bigl(1-117t+3969t^2\bigr)^{1/3}=\theta_{-7}(\tau),
\]
with $\theta_{-7}$ the weight-one theta series of $\Q(\sqrt{-7})$, the unique
form in $M_1(\Gamma_0(7),\chi_{-7})$; equivalently
$\thq t/t=\theta_{-7}^2\in M_2(\Gamma_0(7))$. Moreover
\[
\frac{t\,\theta_{-7}(\tau)^3}{1-117t+3969t^2}=-\tfrac19\,\eta(\tau)^3\eta(7\tau)^3,
\]
the weight-three newform \textup{7.3.b.a}. Since
$\lambda^2-117\lambda+3969$ has discriminant $-2187$, $|\lambda_1|=|\lambda_2|=63$
and there is no archimedean Ap\'ery limit; $k=2$. The row is integral only after
the scaling $\lambda=3$.
\item $(n+1)^2u_{n+1}=(72n^2+36n+6)u_n-108(4n-1)(4n-3)u_{n-1}$, $u_0=1$,
\[u_n=1,\,6,\,90,\,1140,\,-5850,\,-1265004,\dots\in\Z ,\]
class $(\tfrac12;\tfrac12)$, $\mathcal I=3$, $\deg\mathcal J=3$, $h=3$: the
surface $I_3\,III\,III\,III$ \textup{(}\#45, $\mathcal J=X^3/Y^3$\textup{)},
projective monodromy the index-three congruence group of level three, the kernel
of $\PSL_2(\Z)\twoheadrightarrow\Z/3$. Uniformisation:
\[
j=\frac{(24t-1)^3}{8t^3},\qquad 3\,\thq t/t=A(t)^2(1-72t+1728t^2)^{1/2},
\]
with $t=-\tfrac12w-6w^2-72w^3-740w^4-\cdots$ a series in $w=q^{1/3}$.
$\lambda^2-72\lambda+1728$ has discriminant $-1728$, $|\lambda_i|=41.57$, no
archimedean limit; $k=2$; integral only after $\lambda=6$.
\end{enumerate}
\end{proposition}

The two behave differently, and the difference is the gauge exponent. At \#30
it is $\tfrac13$ and the two $II$ fibres are order-$3$ elliptic points, so
$R^{1/3}$ is single-valued on $\HH$ and $A\cdot R^{1/3}$ is a genuine holomorphic
modular form; at \#45 it is $\tfrac14$ while the three $III$ fibres are
order-$2$ points, so $R^{1/4}$ is not single-valued and the row's ``$F$'' is a
square root of a weight-two object --- the quadratic twist that
Theorem~\ref{thm:kodaira} allows. In neither case is the Eichler source
$\Phi=D^2(B/A)$ a weight-three modular form on the relevant group: for \#30 its
expansion $-\tfrac19q-\tfrac4{27}q^2-\tfrac{17}{81}q^3+\cdots$ has unbounded
$3$-power denominators and the fit against $M_3(\Gamma_0(7),\chi_{-7})$ fails
from $q^4$ on. Since the characteristic roots are complex conjugate the natural
constant is not an Ap\'ery limit but a \emph{fold constant}, the connection
coefficient at the complex singular point $t_c$; both folds sit at classical CM
points ($j(\tau_*)=1728$, discriminant $-4$, for \#45, as it must be with the
$III$ fibres over $j=1728$; $\tau_*\sim\rho$, discriminant $-3$, for \#30), so
Chowla--Selberg periods are legitimate candidates alongside $L$-values. For \#45
the identification succeeds.

\begin{proposition}[a Chowla--Selberg period]\label{prop:CS}
\Numerical{$\ge245$ digits; agreement $1.0\times10^{-249}$} For the row \#45,
\[
\operatorname{Im}\xi=-\frac{\Gamma(1/4)^4}{2^{7/2}\,3^{11/4}\,\pi}
=-\frac{\varpi^2}{\sqrt2\,3^{11/4}},
\]
$\varpi=\Gamma(1/4)^2/(2\sqrt{2\pi})$ the lemniscate constant, the
Chowla--Selberg period of discriminant $-4$; equivalently
$4\cdot3^{11}(\operatorname{Im}\xi)^4=\varpi^8$, and \texttt{algdep} returns
$2^{14}3^{11}x^4-1$ for $\operatorname{Im}\xi/(\Gamma(1/4)^4\pi^{-1})$ at degree
$4$, height $2\times10^9$.
\end{proposition}

This is the first constant in this census that is a Chowla--Selberg period rather
than a critical $L$-value --- a direct consequence of the source not being
modular. The corresponding constants for \#30 resist identification against every
basis tried (Chowla--Selberg monomials, Dirichlet and newform $L$-values,
\texttt{algdep} with algebraic coefficients to degree $18$). \Open{}

Finally, the class $(\tfrac12;\tfrac12)$ has Fuchsian signature
$(0;2,2,2;\text{one cusp})$, and \emph{two} groups of that signature occur among
integral rows: the index-three subgroup of $\PSL_2(\Z)$ above (row 6) and
$\Gamma_0(5)+5$, of the same covolume but not conjugate into $\PSL_2(\Z)$,
carrying the level-$5$ Fricke row (row 10). The $\mathcal J$-map test is exactly
what separates them.

\subsection{The classification theorem}

Theorem~\ref{thmC} in full; parts (i)--(ii) are Theorems~\ref{thmA}
and~\ref{thmB} and are not repeated.

\begin{theorem}[classification of second-order rows]\label{thm:herfclass}
Let \eqref{eq:R2gen} have $u_n\in\Z$ for all $n$ \textup{(}after a $\lambda^n$
rescaling\textup{)}, $\deg P,\deg Q\le2$, $d\neq0$, $a^2\neq4d$, $Q(n)\neq0$ for
$n\ge1$.
\begin{enumerate}[label=\textup{(\roman*)},leftmargin=2.6em,start=3]
\item \VerifiedR{all $26$ Kodaira-admissible classes, $1\le A\le3000$,
$|B|\le150$, $|C|\le\min(4000,12000/M^2)$, integrality prime-tested to $n=30$ and
re-verified exactly to $n=300$} In the classes attached to \emph{rigid}
configurations the complete list of primitive, Casoratian-non-degenerate rows is
the $28$ rows of Table~\ref{tab:rows}. Exactly ten have an elliptic-surface local
system: Zagier's six ($\deg\mathcal J=12$; Beauville's $I_1I_2I_3I_6$ three
times, $I_1I_1I_1I_9$, $I_1I_1I_5I_5$, $I_1I_1I_2I_8$), the two root rows of
Theorem~\ref{thm:whichroots} ($\deg\mathcal J=6$), and the two new rows of
Proposition~\ref{prop:newrows} ($\deg\mathcal J=8$ and $3$). The other eighteen
have projective monodromy commensurable with, but not conjugate into,
$\PSL_2(\Z)$.
\item \VerifiedR{same box} In the classes attached to Herfurtner's
\emph{one-parameter} families the integral rows form infinite families; the
cleanest is $(\rho;\delta_\infty)=(0;1)$, $P(n)=\tfrac A4(2n+1)^2$,
$Q(n)=C(2n-1)(2n+1)$, integral for $A\in4\Z$ and $C$ in an arithmetic
progression --- one parameter after $t\mapsto t/\mu$, matching the one modulus of
Herfurtner's $I_1I_1I_1I_3^*$ / $I_1I_1I_2I_2^*$ families.
\end{enumerate}
\end{theorem}

\begin{corollary}[the only positive scores]\label{cor:H6}
\VerifiedR{complete pass over all $1\,270\,065$ scan hits: a positive score needs
$|\lambda_2|<e^{-k}\le e^{-1}$, a condition on $(A,D)$ alone}
Of the $69\,434$ hits in Casoratian-non-degenerate classes, $1\,374$ have
$|\lambda_2|<1$ with real characteristic roots, and \emph{exactly two} have a
positive score with $k\ge2$:
\[
\text{Ap\'ery's }\zeta(2)\text{ row }(11,3,-1)\text{ at }+0.40606,\qquad
\text{Beukers' row }(136,10,4)\text{ at }+0.13920 .
\]
Every other positive score in the scan has $k=1$ and belongs to the classical
Legendre--Pad\'e $\log$ family (top of that ranking: $A=600$, score $+2.5115$,
$k=1$).
\end{corollary}

\begin{remark}[what the scan does not settle]
The eleven Herfurtner one-parameter families with a moving $\mathcal J$ are not
excluded by Theorem~\ref{thm:finite} --- their cross-ratio is free --- only by the
scan; and a box is a box. \Open{} The proved part of the classification is
entirely local (Lemma~\ref{lem:exps},
Theorems~\ref{thm:norm},~\ref{thm:kodaira},~\ref{thm:finite},
Proposition~\ref{prop:K1}); the global input is Herfurtner's theorem, and the
finiteness of the answer is the finiteness of his list together with the
rationality of $a^2/d$.
\end{remark}

\subsection{What an unconditional classification would need}\label{sec:BD}

Theorem~\ref{thmA} is conditional on the geometric hypothesis;
Theorem~\ref{thm:herfclass} replaces it by a box. Neither is the statement one
wants, which is Zagier's conjecture as quoted in \S\ref{sec:intro}: here his
$P(t)$ is our $R(t)=1-at+dt^2$, so ``non-degenerate'' is exactly $d\ne0$,
$a^2\ne4d$. This is the rank-two, four-singular-point instance of the
Bombieri--Dwork expectation \cite{Bombieri1981} that a globally nilpotent
Fuchsian operator is of geometric origin: an integral row's solution is a
$G$-function, so by the Chudnovsky theorem \cite{Chudnovsky1985} its minimal
operator is globally nilpotent, and what is missing is the passage from that to a
family of varieties. The converse reduction is proved here.

\begin{proposition}\label{prop:reduction}
\Proved{} \textup{(}modulo the classification inputs already used\textup{)}
Suppose an integral row \eqref{eq:R2gen} with $d\ne0$, $a^2\ne4d$ has a modular
parametrisation: $t$ a Hauptmodul of a genus-zero $\Gamma$ commensurable with
$\PSL_2(\Z)$, the solution the $t$-expansion of a weight-one form on $\Gamma$.
Then by Lemma~\ref{lem:singpts} the singular points are the special points of
$\Gamma$, so $\Gamma$ has exactly four of them; by Theorem~\ref{thm:N2} its
covolume index is $\le12$ and, if $\Gamma\subseteq\PSL_2(\Z)$, it is one of $28$
enumerated groups; and Theorems~\ref{thmA} and \ref{thm:N3} then place the row on
the list. Hence \emph{Zagier's conjecture for this class is equivalent to the
statement that the projective monodromy group of an integral three-term row is
commensurable with $\PSL_2(\Z)$.}
\end{proposition}

Why that is hard was set out in \S\ref{sec:intro}: rank two with four singular
points is not rigid in Katz's sense \cite{Katz1996}, the accessory parameter
being the modulus; the implication ``globally nilpotent $\Rightarrow$
hypergeometric pullback'' may not be used as a shortcut, Krammer's Lam\'e
equation \cite{Krammer1996,MovasatiReiter} being the standing warning; and the
strongest partial result, Beukers' $p$-adic analysis of Dwork's accessory
parameter problem in exactly this configuration \cite{Beukers2002}, does not
close it.

What the computations supply instead is a census, and it is consistent: over
every box searched, every integral row is (a) one of the ten elliptic-surface
rows of Theorem~\ref{thm:herfclass}(iii); or (b) a row on an Atkin--Lehner
quotient --- Beukers'/$\sqrt{\text{Ap\'ery}}$, $\sqrt T$, $\sqrt{\text{Domb}}$,
$\sqrt{\mathrm{AZ}(7,3,81)}$, $\sqrt{s_7}$, the level-$5$ Fricke row, six further
new rows and the $A=0$ rows, all still \emph{modular} in Zagier's sense, the
universal curve simply failing to descend to their base; or (c) a member of one
of the two classical Pad\'e families --- the Legendre--Pad\'e $\log$ family, which
occupies the classes attached to Herfurtner's one-parameter $I_m^*/I_0^*$
families, and the $\arctan$ family, which Theorem~\ref{thmB}(ii) excludes from
the elliptic-surface world and whose canonical coordinate is not even integral
($\operatorname{lcm}$ of denominators over $n\le58$ equal to
$2^{11}5^67^411^213^2\cdot17\cdot19\cdot23$, new primes entering indefinitely, so
$t(q)$ is not algebraic by Eisenstein's theorem \VerifiedR{$n\le58$}). No row in
any box is unaccounted for; but Proposition~\ref{prop:reduction} is where the
argument stops. \Open{}

% ===========================================================================
\section{Cusp moves and rigidity}\label{sec:cusp}
% ===========================================================================

Table~\ref{tab:placements} shows three of Zagier's rows on one surface. That is not
a coincidence of cross-ratios but a group action, and the action has a sharp
arithmetic consequence.

\begin{theorem}[the general cusp move]\label{thm:move}
\Proved{} \textup{(}symbolically\textup{)} Let $y$ solve $Ly=0$ for a row
$(P,Q)$ with characteristic roots $\lambda,\mu$ and $Q$-roots $r_1,r_2$. Put
\[
s=\frac{t}{1-\lambda t}\quad\Bigl(\text{so }t=\frac{s}{1+\lambda s}\Bigr),
\qquad z(s)=(1-\lambda t)^{\alpha}y(t).
\]
Then $z$ is again the analytic solution of a three-term row with coefficients of
degree $\le2$ \emph{if and only if}
\[
(\alpha-1+r_1)(\alpha-1+r_2)=0,\qquad\text{i.e.}\qquad \alpha\in\{1-r_1,1-r_2\},
\]
i.e.\ exactly when $\alpha$ is one of the two local exponents of $L$ at
$t=\infty$. With $\alpha=1-r_1$ the new row is
\[
\begin{aligned}
P^{\sharp}(n)&=(\mu-2\lambda)n^2+(p_1-3\lambda+2\lambda r_1)n
   +(p_0-\lambda+\lambda r_1),\\
Q^{\sharp}(n)&=\lambda(\lambda-\mu)n^2
   +\lambda(\lambda+\mu-p_1-2\lambda r_1+\mu r_1-\mu r_2)n\\
&\qquad+\lambda(\lambda r_1^2-\lambda r_1+\mu r_1r_2-\mu r_1+p_1r_1),
\end{aligned}
\]
with characteristic roots $-\lambda$ and $\mu-\lambda$, and the two operators are
related by $L^{\sharp}z=(1-\lambda t)^{1-\alpha}\,Ly$.
\end{theorem}

\begin{proof}[Proof sketch]
Substituting $t=s/(1+\lambda s)$ and $y=(1+\lambda s)^{\alpha}z$ into $Ly$ and
dividing by $(1+\lambda s)^{\alpha-1}$ makes the coefficients of $z''$ and $z'$
polynomial of degree $2$ unconditionally, while the coefficient of $z$ is
polynomial with remainder
\[
\frac{\alpha^2\lambda\mu-2\alpha\lambda\mu-\alpha q_1+\lambda\mu+q_0+q_1}{\lambda^2}
=\frac{d}{\lambda^2}(\alpha-1+r_1)(\alpha-1+r_2)
\]
after $q_1=-d(r_1+r_2)$, $q_0=dr_1r_2$. Setting this to zero is the stated
condition, and the boxed coefficients are then read off. The computation is a
finite identity in the Weyl algebra and was carried out symbolically.
\end{proof}

The reading is transparent: the move sends $t=\infty$ to $s=-1/\lambda$, where a
three-term row needs local exponents $(0,\rho^\sharp)$; the gauge
$(1-\lambda t)^{\alpha}$ shifts the exponents there by $-\alpha$, and the two
admissible choices are the two exponents $1-r_i$. So the move is the exponent
bookkeeping of Theorem~\ref{thm:norm} turned into a group action. In Zagier's
class ($Q=dn^2$, $r_1=r_2=0$, $\alpha=1$) it specialises to
$(a,b,d)\mapsto(\mu-2\lambda,\ b-\lambda,\ \lambda^2-\lambda\mu)$ and, on
coefficients, to
$u^\sharp_n=\sum_m\binom{n+\alpha-1}{n-m}(-\lambda)^{n-m}u_m$, the signed
binomial transform; for $\alpha=1$, $\lambda\in\Z$ integrality is automatic,
while for $\alpha\notin\Z$ the binomials carry denominators at the primes
dividing $\operatorname{den}(\alpha)$ and a rescaling $u_n\mapsto c^nu_n$ with
$c\in\{1,2,4,9\}$ restores it in every case measured.

\begin{theorem}[the companion]\label{thm:companion}
\Proved{}; \VerifiedR{all $22$ available (row, $\lambda$, $\alpha$) cases, exactly
to $n=40$} In the situation of Theorem~\ref{thm:move}, let $\widehat B$ solve
$L\widehat B=t(1-\lambda t)^{-\alpha}$ with $\widehat b_0=0$. Then the companion
of the moved row is $B^{\sharp}=(1-\lambda t)^{\alpha}\widehat B$. In particular
the naive guess $B^{\sharp}=(1-\lambda t)^{\alpha}B$ is \emph{false}: what is
replaced is the inhomogeneity, not the companion.
\end{theorem}

\begin{proof}
By Theorem~\ref{thm:move}, $L^{\sharp}[(1-\lambda t)^{\alpha}X]
=(1-\lambda t)^{1-\alpha}LX$ for any $X$. With $X=\widehat B$,
$L^{\sharp}B^{\sharp}=(1-\lambda t)^{1-\alpha}\cdot t(1-\lambda t)^{-\alpha}
=t/(1-\lambda t)=s$, which is the defining inhomogeneity of the moved row's
companion; and $B^\sharp=s+O(s^2)$. Uniqueness of the analytic solution with
$b_0=0$ finishes.
\end{proof}

Consequently the Ap\'ery limit is \emph{not} preserved by a cusp move: two
things change at once, the dominant singularity and the second solution. In
modular terms, with $\Phi=F\,Dt$ the weight-three source and
$\Theta=D^{-2}\Phi$, $B^\sharp=F^\sharp D^{-2}(\Phi/(1-\lambda t))$;
multiplication by the modular unit $(1-\lambda t)^{-1}$ is not Hecke-equivariant,
so the source can jump between Hecke eigenspaces.

\begin{theorem}[orbits and rigidity]\label{thm:orbits}
\Proved{} In the coordinate $v=1/t$ the four singular points are $v=\infty$
\textup{(}MUM\textup{)}, $0$, $\lambda$, $\mu$, and the cusp move by $\lambda$ is
the translation $v\mapsto v-\lambda$. Hence the orbit consists of exactly three
\emph{placements} --- which of $\{0,\lambda,\mu\}$ goes to $v=0$ --- each carrying
at most four gauges, so at most $12$ rows modulo $u_n\mapsto c^nu_n$; and the
placement at $v=v_j$ has $|\lambda_2|=\min_{i\ne j}|v_i-v_j|$. Moreover
\textup{(}Theorem~\ref{thmB}(iii)\textup{)}, if $a^2-4d$ is a square then all
three gaps $|\lambda|,|\mu|,|\lambda-\mu|$ are positive integers, so
$|\lambda_2|\ge1$ and $\operatorname{score}\le-k\le-1$ at every placement.
\end{theorem}

\begin{corollary}\label{cor:irrational-roots}
\Proved{} No rational cusp-move orbit contains a positive-score row; conversely a
positive-score row has irrational characteristic roots, so its only other
equally-good placement is its Galois conjugate over $K=\Q(\sqrt{a^2-4d})$.
\end{corollary}

Measured orbit sizes in the corpus of $49$ distinct rows are $2,3,8,12$; the
$\mathbf A,\mathbf C,\mathbf F$ orbits coincide, all members are integral to
$n=200$ after a rescaling $c\in\{1,2,4,9\}$, $k\in\{1,2\}$ throughout and sharp,
the best score in a rational orbit is exactly $-2$, and the only corpus rows with
$|\lambda_2|<1$ --- Ap\'ery's $\zeta(2)$ row $(0.09017)$, Beukers' $(0.11775)$,
$\sqrt T$ $(0.68629)$, $\Gamma_0(5)+5$ $(0.72136)$ --- all have irrational roots.
\VerifiedR{$49$ rows, $8$ orbits, $n\le200$} The conjugate placements exist over
$K$ with the \emph{same} score to the last digit but approximate a different real
number: $0.32781709442\ldots$ against $\zeta(2)/5=0.32898681337\ldots$, and
$0.10021505333\ldots$ against $0.10018744923\ldots$ \Numerical{$\ge231$ digits}

\subsection{One surface, two extension classes}

With $t_{\mathbf C}=\eta_1^4\eta_6^8/(\eta_2^8\eta_3^4)$,
$F_{\mathbf C}=\eta_2^6\eta_3/(\eta_1^3\eta_6^2)$ and $\mathbf F'$ denoting
$\mathbf F$ in the level-$6$ normalisation
$u_n^{\mathbf F'}=(-1)^nu_n^{\mathbf F}$,
\[
t_{\mathbf A}=\frac{t_{\mathbf C}}{1-t_{\mathbf C}},\quad
t_{\mathbf F'}=\frac{t_{\mathbf C}}{1-9t_{\mathbf C}},\quad
F_{\mathbf A}=(1-t_{\mathbf C})F_{\mathbf C},\quad
F_{\mathbf F'}=(1-9t_{\mathbf C})F_{\mathbf C}
\]
\VerifiedR{exact to $O(q^{241})$}. The cusps of $\Gamma_0(6)$ are
$\infty,1/2,1/3,0$ of widths $1,3,2,6$, carrying $I_1,I_3,I_2,I_6$;
$t_{\mathbf C}$ vanishes at $\infty$, has its pole at $1/2$ and takes the values
$1$, $\tfrac19$ at $1/3$, $0$, which is why the three substitutions put
$I_3,I_2,I_6$ at $t=\infty$. The relating map is \emph{not} an Atkin--Lehner
involution: $\{1,W_2,W_3,W_6\}$ acts simply transitively on the cusps, so no
nontrivial element fixes the MUM cusp.

The surface fixes the pure local system; the choice of MUM cusp fixes which
Eisenstein class degenerates there. With $\mathcal S=E_{3,\chi_{-3},\mathbf 1}$,
$\mathcal T=E_{3,\mathbf 1,\chi_{-3}}$ and $(V_kf)(\tau)=f(k\tau)$, one has
$\Phi_{\mathbf C}=(1-8V_2)\mathcal S$, $\Phi_{\mathbf F'}=(1+V_2)\mathcal S$,
$\Phi_{\mathbf A}=(1-V_2)\mathcal T$ \VerifiedR{exact to $O(q^{260})$}. The four
series $\mathcal S,V_2\mathcal S,\mathcal T,V_2\mathcal T$ are independent, so
$\{\mathcal S,V_2\mathcal S\}$ and $\{\mathcal T,V_2\mathcal T\}$ span two
different planes of $\mathcal M_3(\Gamma_0(6),\chi_{-3})$; the cusp move
multiplies $\Phi$ by the modular unit $(1-\lambda t)^{-1}$ and is not
Hecke-equivariant, landing in the $\mathcal S$-plane for $\lambda=9$ and in the
$\mathcal T$-plane for $\lambda=1$. Hence \emph{one surface, two extension
classes}: $\xi^{\mathbf C}=\tfrac12L(2,\chi_{-3})$,
$\xi^{\mathbf F'}=-\tfrac58L(2,\chi_{-3})$, so
$\xi^{\mathbf F'}/\xi^{\mathbf C}=-\tfrac54$ exactly, while
$\xi^{\mathbf A}=\tfrac14\zeta(2)$ and $\xi^{\mathbf A}/\xi^{\mathbf C}$ has no
rational approximation of height $\le10^{12}$ beyond the random level.
\Numerical{$\ge80$ digits}

\begin{corollary}[the $3$-adic relation]\label{cor:3adic}
Inside the $\mathcal S$-plane,
$\Theta_{\mathbf F'}=-\tfrac54\Theta_{\mathbf C}+\tfrac94\Theta_{\Psi_0}$ with
$\Psi_0=(1-4V_2)\mathcal S$ the class of vanishing Mellin factor
\textup{(}\VerifiedR{$O(q^{260})$}\textup{)}; multiplying by
$F_{\mathbf F'}=(1-9t_{\mathbf C})F_{\mathbf C}$ gives
\[
B^{\mathbf F'}+\tfrac54\xi A^{\mathbf F'}
=(1-9t_{\mathbf C})\Bigl(-\tfrac54R^{\mathbf C}_\xi+\tfrac94\Xi\Bigr),
\qquad t_{\mathbf C}=\frac{t_{\mathbf F'}}{1+9t_{\mathbf F'}} .
\]
Since $x\mapsto x/(1+9x)$ is an isometric automorphism of the disc
$\{|x|_3\le3^{\delta}\}$, $0<\delta<2$, and $1\pm9x$ are units of norm $1$ there,
$\xi_3^{\mathbf F}=\tfrac54\xi_3^{\mathbf C}$ --- conditionally on
overconvergence of $R^{\mathbf C}_\xi$ and $\Xi$ on such a disc and on
$v_3(a_n^{\mathbf F})=O(\log n)$.
\Numerical{$v_3(4\xi_3^{\mathbf F}-5\xi_3^{\mathbf C})=1586$ at $N=800$, against
Cauchy precision $1581$}
\end{corollary}

The move transports nothing to $\mathbf A$, and cannot: $\Theta_{\mathbf C},
\Theta_{\Psi_0},\Theta_{\mathbf A}$ have rank $3$, so no such identity exists;
$\mathbf A$ has $\sigma_3=v_3(d)=0$ and no $3$-adic Ap\'ery limit at all; and at
its own alignment prime $\xi_2^{\mathbf A}=0$. The failure is structural, not a
gap.

\subsection{The cusp move can change the source's Hecke type}
The $\sqrt{\text{Domb}}$ orbit (size $8$) contains two members with genuinely
mixed limits: $(20,10,64)$ with
$\xi=\tfrac1{16}(15L(2,\chi_{-3})-6\zeta(2))$ and $(20,4,64)$ with
$\xi=\tfrac1{32}(6\zeta(2)+15L(2,\chi_{-3}))$, both integral to $n=200$ with
$k=2$ sharp, $\lambda_1=16$, $\lambda_2=4$.
\Numerical{agreement $2.8\times10^{-124}$ resp.\ $2.7\times10^{-118}$ at $n=200$
against predicted truncation $3.9\times10^{-121}$} Their parent carries the
\emph{cuspidal} value $L(f_{12},2)$: the cusp move has pushed the weight-three
source out of the cuspidal part into the Eisenstein part of
$M_3(\Gamma_0(12),\chi_{-3})$, which Theorem~\ref{thm:companion} explains. These
are the first rows in this census whose single Ap\'ery limit is a nontrivial
$\Q$-combination of $\zeta(2)$ and $L(2,\chi_{-3})$, with independent coefficient
vectors $(-6,15)$ and $(6,15)$ --- the shape a Nesterenko argument for
$\dim\langle1,\zeta(2),L(2,\chi_{-3})\rangle$ wants. It does not get near: here
$\sigma=2+\log16$ and $\delta=-2-\log4$, a deficit of $3.386$ nats per step to
$\dim\ge2$, and from the better pair $(\mathbf A,\mathbf C)$ the deficit to
$\dim\ge3$ --- the theorem of Calegari--Dimitrov--Tang \cite{CDT2024} --- is
$10.394$ nats per step.

% ===========================================================================
\section{Non-congruence hosts}\label{sec:noncong}
% ===========================================================================

The classification above is organised by the surface. There is a second
organisation, by the group, and it produces the same finiteness from
Gauss--Bonnet rather than from Riemann--Hurwitz.

Let $\Gamma$ be a genus-zero Fuchsian group commensurable with $\PSL_2(\Z)$ with
a distinguished cusp, $t$ a Hauptmodul normalised $t=q_h+O(q_h^2)$ with integral
coefficients, and $F$ a weight-two form on $\Gamma$ with $F=1+O(q)$ and divisor
supported on the cusps (an eta quotient, say). Put $g=\sqrt F$ and
$a_n=\lambda^n[t^n]g$. Then $g$ is a weight-one form on an index-$\le2$ subgroup,
$a_n$ satisfies a second-order row, and $4^n[t^n]\sqrt G\in\Z$ for any
$G\in1+x\Z[[x]]$ --- sharply, since $4^k\binom{1/2}k=(-1)^{k-1}2C_{k-1}$ with
$C_j$ Catalan --- so $\lambda=\max(1,2^{2-e})\in\{1,2,4\}$ where
$e=\min_{n\ge1}v_2([t^n]F)$.

\begin{lemma}\label{lem:singpts}
\Proved{} The finite singular points of the weight-one operator in the $t$-line
are exactly the values $t(P)$ for $P$ a cusp or elliptic point of $\Gamma$ with
$t(P)\notin\{0,\infty\}$. They do \emph{not} depend on $F$.
\end{lemma}

\begin{proof}
At a non-elliptic $P\in\HH$, $t$ is a local isomorphism and $g$ is holomorphic
and non-vanishing, so both solutions are holomorphic and independent and $t(P)$
is ordinary. At an elliptic point of order $e$ the local coordinate on
$X_\Gamma$ is $w\sim(\tau-\tau_P)^e$ and $t-t(P)\sim w$, so the local monodromy
is a rotation by $2\pi/e$ and the exponent difference is $1/e$. At a cusp with
$t$ finite the second solution carries $\log(t-t(c))$. Zeros or poles of $F$
change the exponents, not the locations.
\end{proof}

\begin{theorem}[the score splits]\label{thm:N1}
\Proved{} With $k=2$ as above and $|t_1|\le|t_2|\le\cdots$ the singular values,
\[
\operatorname{score}=\log(1/|\lambda_2|)-k=\log|t_2|-2-\log\lambda .
\]
The archimedean factor $|t_2|$ depends only on $(\Gamma,t)$ --- not on $F$ --- and
$\lambda\in\{1,2,4\}$ depends only on the $2$-adic structure of $F$. Hence a
positive score needs $|t_2|>e^2\lambda$, i.e.\ $|t_2|>7.389$, $14.778$ or
$29.556$ according as $\lambda=1,2,4$.
\end{theorem}

By Calegari--Dimitrov--Tang \cite{CDT2021}, a modular form on a finite-index
subgroup with algebraic Fourier coefficients and \emph{bounded} denominators is a
form on a congruence subgroup. So $\lambda>1$ is a certificate of
non-congruence; the converse direction, and the value of $\lambda$, come from the
elementary binomial estimate above, not from their theorem.

\begin{theorem}[four special points; finiteness]\label{thm:N2}
\Proved{} \textup{(}Gauss--Bonnet\textup{)}$+$\VerifiedR{exhaustive enumeration}
A three-term row forces $\Gamma$ to have exactly four special points (cusps plus
elliptic points), whence its covolume index in $\PSL_2(\Z)$ is
\[
\mu=6\Bigl(2-\sum_{i}\frac1{e_i}\Bigr)\ \le\ 12 ,
\]
with equality exactly in the four-cusp case. Inside $\PSL_2(\Z)$ there are, up to
conjugacy, exactly $28$ genus-zero subgroups with four special points: $16$
congruence and $12$ non-congruence. The non-congruence ones have index $7$, $9$,
$10$ and signatures $(0;2,3;2)$, $(0;2;3)$, $(0;3;3)$, four groups each. The
index-$12$ four-cusp case returns exactly six groups, with cusp widths
$[1,1,5,5]$, $[1,2,3,6]$, $[1,1,2,8]$, $[1,1,1,9]$, $[2,2,4,4]$, $[3,3,3,3]$,
\emph{all congruence} --- Beauville's six, recovered from scratch. In particular
no non-congruence group has four cusps.
\end{theorem}

The enumeration is by transitive pairs $(s,t)\in S_n$ with $s^2=t^3=1$,
canonicalised by breadth-first relabelling, cusps read off as cycles of $st$, and
congruence decided by Wohlfahrt's criterion.

\begin{theorem}[unequal exponents kill the score]\label{thm:N3}
\Proved{} \textup{(}modulo the normalisation $P,Q,R\in\Z[t]$, $P(0)=1$, true for
every row here\textup{)} If the two finite singular points carry different
exponent differences $1/e_1\ne1/e_2$ then $t_1,t_2\in\Q$, hence
$\lambda_1,\lambda_2\in\Z$, so $|\lambda_2|\ge1$ and score $\le-k<0$.
\end{theorem}

This is Corollary~\ref{cor:galois} read on the group side, and it settles the
non-congruence host question. Index $7$ has signature
$(0;2,3;2\ \text{cusps})$: whichever of the three non-MUM special points goes to
$t=\infty$, a \emph{mixed} pair remains at $t_1,t_2$, so all four index-$7$
groups are excluded outright; for index $9$ and $10$ the only non-mixed option is
to send the elliptic point to $\infty$, leaving two cusps at $t_1,t_2$ --- the
Zagier class, covered exhaustively by \S\ref{sec:scan}. So \emph{a
non-congruence host group contributes nothing beyond the Zagier class: the
interesting non-congruence-ness is not in the group of $t$ but in the multiplier
system of $g=\sqrt F$.}

That is exactly Beukers' situation. His $t=(\eta_1\eta_6/\eta_2\eta_3)^{12}$
lives on the \emph{congruence} group $\Gamma_0(6)+6$; his $\sqrt E$ does not.
With $F=\eta_2^7\eta_3^7/(\eta_1^5\eta_6^5)$ Ap\'ery's weight-two form
($F|_2W_6=-F$, divisor $(1/2)+(1/3)$),
$a_n=4^n[t^n]\sqrt F=1,10,534,40900,3672550,\dots$,
$\lambda_{1,2}=4(17\pm12\sqrt2)$, $k=2$, $\lambda=4$, and the score is
$+0.13920$: the criterion $\lambda e^k<|t_2|$ reads $4e^2<(\sqrt2+1)^4$, i.e.\
$29.5562<33.9706$, a margin of $1.1494$. Level $6$ wins because $t\circ W_6=t$
with $t_1t_2=1$ exactly, so $|t_2|=1/|t_1|$ is as large as the fastest-growing
row allows. This is Beukers' Theorem~3 of 1987 \cite{Beukers1987}: the Ap\'ery
limit $\xi=0.100187449229339406\ldots=L(\Psi,2)$, for the weight-three
half-integral eta quotient
$\Psi=\tfrac14(\eta_1\eta_6)^{9/2}(\eta_2\eta_3)^{-3/2}$, is irrational with
$\mu(\xi)\le50.654$.

\begin{theorem}[uniqueness]\label{thm:unique}
\VerifiedR{two independent exhaustive scans; boxes stated below}
Beukers' row is the only positive-score row whose form is non-congruence, and the
only positive-score rows at all with $k\ge2$ are it, Ap\'ery's $\zeta(2)$ row, and
members of two infinite classical Pad\'e families.
\begin{itemize}[leftmargin=1.6em]
\item \emph{Ansatz-free recurrence scan.} Integer $(\alpha,\gamma,\delta,\zeta)$
in the five classes $e\in\{2,3,4,6,\infty\}$ and in a free-exponent variant, at
$\alpha\le3000$, $|\delta|\le150$--$400$, $|\gamma|\le150$--$1000$,
$|\zeta|\le300$--$2000$, with the Casoratian filter and measured $k$, plus a
per-$\rho$ sweep over $\rho=p/q$ with $q\le6$, $0<\rho\le3$. Every positive score
with $k\ge2$ is Ap\'ery's row, Beukers' row, a member of the Legendre--Pad\'e
$\log$ family ($\delta=16$, $\rho=0$, $\xi=-\tfrac1{12}\log\frac{x+1}{x-1}$ at
$\alpha=8x$), or a member of the $\arctan$ family of
Corollary~\ref{cor:excluded}(b). In the square-root class $\rho=\tfrac12$,
Beukers' row is the \emph{unique} positive score at $\alpha\le3000$.
\item \emph{Modular scan.} All eta-quotient pairs $(t,F)$ on $\Gamma_0(N)$,
$N\le60$, with $t$ of weight $0$, $\ord_\infty t=1$, divisor degree $\le4$, and
$F$ holomorphic of weight $2$ with $\ord_\infty F=0$: $3\,036\,854$ pairs. Of
these $739$ admit a recurrence of order $\le4$ and degree $\le3$ ($485$ with
$\lambda=1$, $127$ with $\lambda=2$, $127$ with $\lambda=4$), in $104$ distinct
invariant classes. \emph{Exactly two} have $|\lambda_2|<1$: Beukers'
($+0.1392$, $\lambda=4$, non-congruence) and the level-$8$ $T$ row ($-1.6235$,
$\lambda=1$, congruence). Everything else has $|\lambda_2|\ge1$; the next classes
by $|\lambda_2|$ are $1$ exactly, then $1.0635$, $1.0718$, $1.6180$.
\end{itemize}
\end{theorem}

The modular scan has two documented blind spots --- both $t$ and $F$ are required
to be eta quotients --- so the ansatz-free scan is the load-bearing one, and it is
what turns up the runner-up.

\begin{proposition}[the level-$5$ Fricke row]\label{prop:level5}
\VerifiedR{$a_n\in\Z$ to $n=202$; $k$ to $n=200$; period to $260$ digits}
Let $u=(\eta_5/\eta_1)^6$, a Hauptmodul of $\Gamma_0(5)$,
$t=u/(1+22u+125u^2)=q-16q^2+122q^3-\cdots$ the resulting Hauptmodul of
$\Gamma_0(5)+5$, and $F=E_{2,5}=\tfrac14(5E_2(5\tau)-E_2(\tau))$. Then
$g=\sqrt F$ is weight one on an index-two subgroup, $\lambda=2$ is minimal,
$a_n=2^n[t^n]\sqrt F=1,6,210,10500,615510,\dots$ satisfies
$(n+1)^2a_{n+1}=2(44n^2+22n+3)a_n+4(4n-1)(4n-3)a_{n-1}$ with
$\lambda_{1,2}=44\pm20\sqrt5=\pm8\varphi^{\pm5}$, $k=2$ sharp, non-vanishing
Casoratian, and $\xi=0.164306701064342158629932\ldots$ Since $\lambda=2>1$ and
the source has unbounded $2$-power denominators, $g$ and $\Psi$ live on a
non-congruence group \cite{CDT2021}: a genuine second instance of Beukers'
construction, at the other perfect level. But $|t_2|=2.7726\ll2e^2=14.778$, so
the score is $-1.6734$ and no irrationality follows.
\end{proposition}

The comparison is exact: at level $6$ the Fricke pairing is $t_1t_2=1$, at level
$5$ it is $t_1t_2=-\tfrac14$, and the same $\lambda_1$-for-$|t_2|$ exchange loses
a factor $4$. \emph{That factor of $4$ is the whole difference between the two
levels.} The period resisted identification against
$1,\pi,\pi^2,\pi^3,\zeta(3),\log2,\log5,\log\varphi,\sqrt5$ and against
$L(f,s)$, $s=1,2,3$, for every weight-three newform of every level $M\le120$
divisible by $5$. \Open{}

One cheap diagnostic separates the Pad\'e families from the modular rows in a
single series computation: \emph{mirror-map integrality} (\S\ref{sec:BD}). For
Ap\'ery's $\zeta(2)$ row the canonical coordinate is
$q-5q^2+15q^3-30q^4+40q^5-26q^6-\cdots$, Beukers' $\Gamma_1(5)$ Hauptmodul
$q\prod(1-q^n)^{5(\frac n5)}$; for Beukers' own row it is
$q-48q^2+1056q^3-\cdots$, i.e.\ $\tfrac14(\eta_1\eta_6/\eta_2\eta_3)^{12}$ in the
rescaled nome; for the $\arctan$ family it is not integral at all.
\VerifiedR{$n\le58$}

% ===========================================================================
\section{Outlook}\label{sec:outlook}
% ===========================================================================

\subsection*{$K3$ and higher, and five singular points}
Proposition~\ref{prop:K1} and Lemma~\ref{lem:twistparity} close the $K3$ worry
for four singular points, but in a way worth restating precisely: an elliptic
$K3$ over $\PP^1_t$ with the row's four singular fibres is \emph{not} excluded
--- in the root-row class the mixed $III/III^*$ configuration with $\sum e=24$ is
a genuine possibility --- but every such surface is a quadratic twist of a
rational one with the same $\mathcal J$, and $\deg\mathcal J$ is pinned by the
degree bound together with the parity, so the $\deg\le12$ search window was never
the constraint and no new row can arise. What is not closed is a systematic pass
through Shimada--Zhang's extremal elliptic $K3$s \cite{ShimadaZhang} against the
$26$ classes. \Open{} Five singular points means $\deg P$ or $\deg Q=3$, i.e.\ a
four-term recurrence or cubic coefficients; there is no analogue of Herfurtner's
theorem in that generality, and the cross-ratio becomes a two-parameter modulus,
so the finiteness argument of \S\ref{sec:herf} does not survive. At the other
end, Schmickler-Hirzebruch's classification of elliptic surfaces over $\PP^1$
with \emph{three} exceptional fibres \cite{SchmicklerHirzebruch} is the
degenerate boundary: the case $d=0$, where $t=\infty$ merges with a finite
singular point and the row becomes hypergeometric --- Zagier's modular solutions
\#11 $(16,4,0)$, \#14 $(27,6,0)$ and \#20 $(64,12,0)$ in the notation of
Table~\ref{tab:placements}.

\subsection*{A gate on the higher-rank constructions}
For a linear-independence argument one wants a $\Q$-linear combination $H$ of
several Ap\'ery limits whose linear form still decays --- ``fold-regularity''.

\begin{lemma}[fold-regularity is codimension $\mu$]\label{lem:foldreg}
\Proved{} Let $\lambda_1$ be the dominant characteristic root, of multiplicity
$\mu$. Near $x_1=1/\lambda_1$ the local solution space carries a Jordan block of
size $\mu$: solutions $y_1,y_1\log,\dots,y_1\log^{\mu-1}$ up to holomorphic
corrections. A solution has radius of convergence $>|x_1|$ iff \emph{all $\mu$}
of its coordinates on that block vanish, whereas a hypothesised $\Q$-linear
relation among the limits makes exactly the \emph{leading} one vanish. Hence one
relation implies fold-regularity iff $\mu=1$.
\end{lemma}

For $\mu\ge2$ the form does not decay at all --- one gets polynomial convergence
with $|h_n|^{1/n}\to|\lambda_1|$ --- so there is nothing for a holonomy bound, or
for Beukers' criterion, to act on. Verified sharply: a row with $\mu=1$ folds and
its form decays at rate $\lambda_2$ exactly, while $\Sym^3$ of Zagier $\mathbf E$
($\mu=3$) and a $\Sym^2$ with $\mu=2$ do not fold at all.
\VerifiedR{$n$ up to $10^6$} Symmetric powers and Cauchy squares of a single row
therefore raise $\mu$ and are the wrong place to look; genuinely independent rows
on one surface (\S\ref{sec:cusp}) are the right one.

\subsection*{What the classification says for irrationality proofs}
The score $\log(1/|\lambda_2|)-k$ is not a property of the sequence. By
Theorem~\ref{thm:N1} it splits as $\log|t_2|-2-\log\lambda$, and neither factor
depends on the accessory parameter or on the period: $|t_2|$ is fixed by the
surface together with the choice of which fibre sits at $\infty$, and $\lambda$
by the $2$-adic structure of the weight-two form. The score is therefore a
\emph{motivic} invariant --- of the geometry plus the placement --- and the
classification computes it once and for all. Combined with
Corollary~\ref{cor:irrational-roots} (a positive score forces irrational
characteristic roots, hence a real quadratic field, hence no second rational
placement), Theorem~\ref{thm:N2} (only finitely many host geometries, all
enumerated) and Corollary~\ref{cor:H6} and Theorem~\ref{thm:unique} (over every
box, exactly two rows clear zero, and both were known), the honest conclusion is
that the second-order world is closed. Its two positive-score rows are not
improvable by rearrangement inside the list: the cusp move produces no new
rational data (Theorem~\ref{thm:orbits}), merging two placements gives
$\mu\le1+\min_i\sigma_i/\delta_i$, i.e.\ the better of the two with no
convex-combination gain, and the conjugate placements approximate different
numbers. A new irrationality theorem of this shape must come from outside the
second-order box --- from a higher-rank host, or from a bound of
Bost--Charles/Calegari--Dimitrov--Tang type applied to a whole function module
rather than to a single linear form.

% ===========================================================================
\section{Reproducibility}\label{sec:repro}
% ===========================================================================

{\small
\sloppy\raggedright
All computations are reproducible from the project notes and scripts
\cite{ProjNotes}; \textsc{pari/gp} work uses exact rational (or
$\Q(y)/(y^2-ay+d)$) arithmetic except where a real or $p$-adic precision is
quoted. Under \texttt{lattice/herfurtner/}: \texttt{01\_exponents.py},
\texttt{07\_classtable.py}, \texttt{04\_herfurtner.py} (exponents, classes,
Herfurtner's table and its cross-ratios); \texttt{02\_hscan.c} with
\texttt{run\_scan.sh}, \texttt{run\_a0.sh} ($\approx75$\,min, twelve cores),
\texttt{08\_report.py}, \texttt{09\_table.py} (the scan and its exact
re-verification to $n=300$); \texttt{05\_jtest.gp}, \texttt{10\_jtest\_all.py},
\texttt{11\_jall.gp} (the $\mathcal J$-test at $\deg\le12$);
\texttt{13\_fastfilter.py}, \texttt{14\_decay.py} (the complete positive-score
pass); \texttt{15\_k3jtest.gp}, \texttt{15\_k3rows.py}, \texttt{15\_k3run.gp}
(the raised window, $\mathtt{DMAX}=\mathtt{HMAX}=24$, $\mathtt{NTERM}=140$,
$\mathtt{GLIM}=4096$, $\mathtt{KPRIME}=2^{61}-1$, an $\mathbf F_p$ filter with no
false negatives and exact certificates); data in \texttt{out/}. Under
\texttt{lattice/cusp\_move/}: \texttt{01\_general\_move.py} (the symbolic proof
of Theorem~\ref{thm:move}) through \texttt{12\_domb\_pairs.gp}. Under
\texttt{lattice/acf\_one\_surface/}: \texttt{01\_cusp\_move.gp} through
\texttt{08\_census\_check.gp}. Under \texttt{lattice/noncongruence\_scan/}:
\texttt{01\_class\_scan.c}, \texttt{04\_score.py}, \texttt{06\_groups.py} (the
$28$ groups, Wohlfahrt test), \texttt{02\_eta\_enum.py}, \texttt{03b\_exact.gp},
\texttt{09\_level5.gp}, \texttt{10\_ident5.gp}, \texttt{12\_source5.gp},
\texttt{14\_mirror234.gp}, \texttt{16\_arctan.py}. Under
\texttt{lattice/sqrt\_apery/}: \texttt{01\_sym2.py}, \texttt{12\_sharp.gp},
\texttt{06\_fold.gp}. Under \texttt{lattice/number\_field/}:
\texttt{nf\_herfurtner\_uniformisation.gp}, \texttt{cfold.gp} with
\texttt{run\_all.gp}, \texttt{run\_v23.gp}, \texttt{run\_t.gp}, and the sweeps
\texttt{hunt3.gp}--\texttt{hunt9.gp} (\texttt{xi\_targets.txt}).
}

% ===========================================================================
\begin{thebibliography}{99}
\footnotesize
\setlength{\itemsep}{0pt plus 0.5pt}

\bibitem{AlmkvistZudilin2006}
G.~Almkvist and W.~Zudilin,
\emph{Differential equations, mirror maps and zeta values},
in: Mirror Symmetry V, AMS/IP Stud. Adv. Math. \textbf{38} (2006), 481--515.
arXiv:math/0402386.

\bibitem{Beauville1982}
A.~Beauville,
\emph{Les familles stables de courbes elliptiques sur $\PP^1$ admettant quatre
fibres singuli\`eres},
C. R. Acad. Sci. Paris S\'er. I Math. \textbf{294} (1982), 657--660.

\bibitem{Beukers1987}
F.~Beukers,
\emph{Irrationality proofs using modular forms},
Journ\'ees arithm\'etiques de Besan\c{c}on, Ast\'erisque \textbf{147--148}
(1987), 271--283.

\bibitem{Beukers2002}
F.~Beukers,
\emph{On Dwork's accessory parameter problem},
Math. Z. \textbf{241} (2002), 425--444.

\bibitem{Bombieri1981}
E.~Bombieri,
\emph{On $G$-functions},
in: Recent Progress in Analytic Number Theory, vol.~2, Academic Press (1981),
1--67.

\bibitem{CDT2021}
F.~Calegari, V.~Dimitrov and Y.~Tang,
\emph{The unbounded denominators conjecture},
preprint (2021), arXiv:2109.09040.

\bibitem{CDT2024}
F.~Calegari, V.~Dimitrov and Y.~Tang,
\emph{The linear independence of $1$, $\zeta(2)$ and $L(2,\chi_{-3})$},
preprint (2024), arXiv:2408.15403.

\bibitem{Cooper2012}
S.~Cooper,
\emph{Sporadic sequences, modular forms and new series for $1/\pi$},
Ramanujan J. \textbf{29} (2012), 163--183.

\bibitem{Chudnovsky1985}
D.~V.~Chudnovsky and G.~V.~Chudnovsky,
\emph{Applications of Pad\'e approximations to the Grothendieck conjecture on
linear differential equations},
in: Number Theory (New York 1983--84), Lecture Notes in Math. \textbf{1135},
Springer (1985), 52--100.

\bibitem{Herfurtner1991}
S.~Herfurtner,
\emph{Elliptic surfaces with four singular fibres},
Math. Ann. \textbf{291} (1991), 319--342.

\bibitem{Katz1996}
N.~M.~Katz,
\emph{Rigid local systems},
Ann. of Math. Studies \textbf{139}, Princeton Univ. Press, 1996.

\bibitem{Kodaira1963}
K.~Kodaira,
\emph{On compact analytic surfaces II, III},
Ann. of Math. \textbf{77} (1963), 563--626; \textbf{78} (1963), 1--40.

\bibitem{Krammer1996}
D.~Krammer,
\emph{An example of an arithmetic Fuchsian group},
J. Reine Angew. Math. \textbf{473} (1996), 69--85.

\bibitem{MirandaBook}
R.~Miranda,
\emph{The basic theory of elliptic surfaces},
ETS Editrice, Pisa, 1989.

\bibitem{MovasatiReiter}
H.~Movasati and S.~Reiter,
\emph{Heun equations coming from geometry},
Bull. Braz. Math. Soc. (N.S.) \textbf{43} (2012), 423--442; arXiv:0902.0760.

\bibitem{SchmicklerHirzebruch}
K.~Schmickler-Hirzebruch,
\emph{Elliptische Fl\"achen \"uber $\PP_1(\mathbb C)$ mit drei Ausnahmefasern
und die hypergeometrische Differentialgleichung},
Schriftenreihe Math. Inst. Univ. M\"unster, 1985.

\bibitem{Verrill}
H.~A.~Verrill,
\emph{Picard--Fuchs equations of some families of elliptic curves},
in: Proceedings on Moonshine and related topics (Montr\'eal 1999), CRM Proc.
Lecture Notes \textbf{30}, Amer. Math. Soc. (2001), 253--268.

\bibitem{ShimadaZhang}
I.~Shimada and D.-Q.~Zhang,
\emph{Classification of extremal elliptic $K3$ surfaces and fundamental groups of
open $K3$ surfaces},
Nagoya Math. J. \textbf{161} (2001), 23--54.

\bibitem{VidunasFilipuk}
R.~Vidunas and G.~Filipuk,
\emph{A classification of coverings yielding Heun-to-hypergeometric reductions},
preprint (2012), arXiv:1204.2730.

\bibitem{Zagier2009}
D.~Zagier,
\emph{Integral solutions of Ap\'ery-like recurrence equations},
in: Groups and Symmetries, CRM Proc. Lecture Notes \textbf{47}, Amer. Math. Soc.
(2009), 349--366.

\bibitem{ProjNotes}
Claude (Opus 5), Claude (Fable) and River,
\emph{Project notes and scripts} (2026):
\texttt{HERFURTNER\_CLASSIFICATION.md} (the exponent dictionary, the cross-ratio
invariant, the scan, the $\mathcal J$-map test),
\texttt{HERFURTNER\_K3\_WINDOW.md} (the Riemann--Hurwitz bound and the
$\deg\le24$ recheck),
\texttt{CUSP\_MOVE\_PROGRAM.md} (the cusp move, orbits, rigidity),
\texttt{ACF\_ONE\_SURFACE.md} ($\mathbf A,\mathbf C,\mathbf F$ on one surface),
\texttt{NONCONGRUENCE\_SCAN.md} (Theorems N1--N3, the $28$ groups, the two
scans),
\texttt{SQRT\_APERY.md} (Beukers' row),
\texttt{NUMBER\_FIELD\_PROGRAM.md} (uniformisation of the two new rows, the
Chowla--Selberg period),
\texttt{HOLONOMY\_LINDEP.md} (fold-regularity);
scripts under \texttt{lattice/}, as itemised in \S\ref{sec:repro}.

\end{thebibliography}

\end{document}
